% =========================================================================
% tikz_style.tex — shared visual system for the Section III method figures.
% Usage: % =========================================================================
% tikz_style.tex — shared visual system for the Section III method figures.
% Usage: % =========================================================================
% tikz_style.tex — shared visual system for the Section III method figures.
% Usage: % =========================================================================
% tikz_style.tex — shared visual system for the Section III method figures.
% Usage: \input{tikz_style} in the preamble of a standalone figure.
%
% -------------------------------------------------------------------------
% DESIGN NOTES — conventions extracted from an accepted TGRS manuscript
% (Gao et al., "BBox-free ... SAR ship instance segmentation", TGRS 2024;
%  local copy: scratchpad/gao_bboxfree_tgrs24.pdf).  Observations taken from
%  its Fig. 2 (overall architecture, p. 5), Fig. 5 (FD module, p. 8) and
%  Fig. 6 (conceptual mask panel, p. 8):
%
%  O1  REGION PARTITION.  The architecture figure is carved into 3-4 named
%      regions by DOTTED rectangles with no fill; the region name is set in
%      bold OUTSIDE the rectangle ("Feature Extractor", "Classification
%      Branch", "Localization Branch", "Center Point Extractor").  Nothing
%      is nested more than two levels deep.
%  O2  COLOR CARRIES MEANING, NOT DECORATION.  Only ~5 hues appear, all as
%      pale fills with black text: light gray = ordinary/standard module,
%      salmon-orange = the module the paper contributes (Feature Decoupling
%      Module, Center Point Extractor), yellow = a distinguished operation
%      (Feature Fusion), light blue / light pink = tensors.  Saturated color
%      is reserved for the few glyphs that must be found first.
%  O3  LEGEND PANEL.  Both Fig. 2 and Fig. 5 carry an explicit legend inside
%      the figure (dotted box top-left in Fig. 2, dashed box right in Fig. 5)
%      that decodes arrow colors and single-letter glyph nodes.  The body of
%      the diagram is kept uncluttered by pushing vocabulary into the legend.
%  O4  TENSOR ICONOGRAPHY.  Feature maps are drawn as thin vertical plates,
%      multi-scale features as 3 offset stacked planes, channel vectors as
%      segmented 3-D bars, the backbone as an isometric slab.  Real SAR
%      chips / heat maps / GT overlays are pasted directly into the flow.
%  O5  TENSOR NAMES SIT BESIDE THE GLYPH, NEVER INSIDE IT, in math italic
%      (F_out, F_mask, F_param, F_H_zip ...), placed above-right of the plate.
%  O6  ARROW DISCIPLINE.  Data flow is a thin gray/black line with a small
%      solid triangular head; routing is orthogonal (Manhattan) with a shared
%      trunk when one source fans out to several targets, so lines never
%      cross.  Semantically special transitions get a THICK COLORED BLOCK
%      arrow instead (2x upsample / 2x downsample / connect), i.e. arrow
%      style is itself a semantic channel.  Long skip paths are routed as
%      rails around the outside of the diagram (Fig. 5 top and bottom).
%  O7  TYPOGRAPHY.  One size for block labels, one smaller size for
%      sub-labels, math only in math mode; no stretched or rotated text
%      except a rare 90-degree axis label.
%  O8  ASPECT RATIO.  Double-column method figures are wide and short:
%      Fig. 2 is 530 x 200 pt (2.65:1), Fig. 5 is 530 x 215 pt (2.5:1).
%      They are laid on a grid with equal gutters between parallel branches.
%  O9  CAPTION DIVISION OF LABOR.  The caption names the parts and decodes
%      abbreviations/colors; it does not re-derive the mechanism, which is
%      left to the body text.  The figure is readable without the caption.
%  O10 CONCEPTUAL PANELS (Fig. 6) use dashed sub-panels with a centered
%      label underneath, a pale blue blob for the true object and a pale red
%      overlay for the error region -- schematic, not photorealistic.
%
% -------------------------------------------------------------------------
% WHAT WE ADOPT / WHERE WE DEVIATE
%  + adopt O1 (dotted named regions), O3 (in-figure legend), O4/O5 (tensor
%    glyphs with the symbol outside), O6 (one flow style + Manhattan
%    routing), O8 (wide-short double-column ratios), O10 (schematic panels).
%  ~ deviate on O2: instead of hue-as-identity we bind hue to ROLE in the
%    contribution -- frozen / training-free operator / trained corrector /
%    data -- because the entire argument of this article is which parts are
%    touched and which are not.
%  ~ deviate on the sans-serif block labels of Fig. 2: we set every figure in
%    the body Times face (as the same paper's Fig. 6 does), so figure text
%    and caption text share one face.
%  + every figure is authored at its EXACT printed width (18.2 cm double
%    column = 516 pt, 8.89 cm single column = 252 pt) and included at
%    scale 1, so \footnotesize is a true 8 pt and \scriptsize a true 7 pt on
%    the page in ALL figures.  Use \figcanvasd / \figcanvasc below.
%
% -------------------------------------------------------------------------
% SEMANTIC PALETTE (Okabe-Ito colorblind-safe base; redundant encoding by
% border weight, dash pattern and badge glyph so it survives grayscale):
%   froz  gray   + lock glyph        frozen, not modified by us  (Stage A)
%   ops   vermil + solid heavy rule  our TRAINING-FREE operator  (rung 1)
%   trn   green  + dashed rule       our OPTIONAL TRAINED corrector (2-3)
%   data  blue   + thin rule         data, tensors, I/O
%   warn  red    + dotted rule, small  failure mode / negative control
% =========================================================================

% Real rasters are embedded on the front faces of the tensor glyphs (see
% \realpanel / \fmslab below).  graphicx must be loaded in the preamble, which
% is where this file is \input.
\usepackage{graphicx}
% Single switch for the raster payload: \def\assetdir{assets_light}
% before \begin{document} swaps every embedded raster for its
% print-sized ~600 dpi copy (see figs/make_light_assets.py).  The
% default is the native dump, 965-1016 dpi at the drawn size.
\providecommand{\assetdir}{assets}
\usetikzlibrary{positioning,arrows.meta,fit,backgrounds,calc,shapes.geometric,
                decorations.pathreplacing,decorations.markings,patterns}

% ---------------- palette ----------------
\definecolor{opsred}{HTML}{D55E00}   % Okabe-Ito vermillion  -- training-free ops
\definecolor{trngreen}{HTML}{009E73} % Okabe-Ito bluish green-- trained correctors
\definecolor{datablue}{HTML}{0072B2} % Okabe-Ito blue        -- data / tensors
\definecolor{frozgray}{HTML}{6E6E6E} % neutral               -- frozen base
\definecolor{flowink}{HTML}{262626}  % near black            -- data-flow arrows
\definecolor{warnred}{HTML}{9E2B2B}  % muted brick           -- failure modes
\definecolor{sarink}{HTML}{2B2B33}   % SAR chip background

% ---------------- exact-width canvases ----------------
% \figcanvasd{H} : double column, 18.2 cm = 516 pt = \textwidth  (scale 1.0)
% \figcanvasc{H} : single column,  8.89 cm = 252 pt = \columnwidth
\newcommand{\figcanvasd}[1]{\useasboundingbox (0,0) rectangle (18.2,#1);}
\newcommand{\figcanvasc}[1]{\useasboundingbox (0,0) rectangle (8.89,#1);}

\tikzset{
  % ---- picture-level defaults: ONE font size system for all five figures ----
  fig/.style={font=\footnotesize, line width=0.65pt,
              >={Stealth[length=1.7mm,width=1.35mm]},
              every node/.append style={align=center, inner sep=2.4pt}},
  % ---- blocks: role -> color + redundant border encoding ----
  % Fill tints were raised in v3: at 4-8% the five roles were separated by an
  % RGB distance of only 6.6-21 (out of 441) and all landed on gray 243-248,
  % i.e. the fill channel was consuming ink and returning nothing.  The border
  % weight / dash pattern / badge remain the PRIMARY, grayscale-safe channel.
  ops/.style   ={draw=opsred,   fill=opsred!16,   line width=0.9pt,
                 rounded corners=1.6pt, inner sep=3pt, align=center},
  trn/.style   ={draw=trngreen, fill=trngreen!15, line width=0.9pt,
                 dash pattern=on 1.7pt off 1.1pt,
                 rounded corners=1.6pt, inner sep=3pt, align=center},
  froz/.style  ={draw=frozgray, fill=black!8,     line width=0.7pt,
                 rounded corners=1.6pt, inner sep=3pt, align=center},
  data/.style  ={draw=datablue, fill=datablue!14, line width=0.55pt,
                 rounded corners=1.6pt, inner sep=3pt, align=center},
  plain/.style ={draw=black!45, fill=white,       line width=0.55pt,
                 rounded corners=1.6pt, inner sep=3pt, align=center},
  % v4: radius 1.4 -> 1.6 pt, the one block style that did not share the
  % system's corner radius; weight 0.5 -> 0.6 pt so the failure chips are not
  % the faintest object on the page while being the ones a reader must find.
  % The dotted dash pattern stays as the grayscale-safe channel.
  warn/.style  ={draw=warnred,  fill=warnred!9,   line width=0.6pt,
                 dash pattern=on 1.1pt off 1.1pt,
                 rounded corners=1.6pt, inner sep=3pt, align=center},
  gate/.style  ={draw=opsred,   fill=opsred!16,   line width=0.9pt,
                 diamond, aspect=2.15, inner sep=1pt, align=center},
  % ---- named region (O1): dotted, unfilled, label set outside ----
  grp/.style   ={draw=black!45, densely dotted, line width=0.7pt,
                 rounded corners=2pt, inner sep=0pt},
  grptitle/.style={font=\footnotesize\bfseries, text=black!80, inner sep=1.5pt},
  % ---- step badges (redundant, grayscale-safe stage identity) ----
  badge/.style ={circle, inner sep=0pt, minimum size=3.3mm, text=white,
                 font=\scriptsize\bfseries},
  badgeops/.style={badge, fill=opsred},
  badgetrn/.style={badge, fill=trngreen},
  % ---- typography (O7): exactly three text sizes in the whole figure set ----
  % v3: 7 pt at 55% black was the worst corner of the design space (minimum
  % size AND minimum contrast) -- sub-labels are now 7 pt at 70% black, and
  % anything carrying a QUANTITY is promoted to 8 pt via \qty / \subk.
  ttl/.style   ={font=\footnotesize\bfseries},
  lbl/.style   ={font=\footnotesize},
  sub/.style   ={font=\scriptsize, text=black!70},
  subk/.style  ={font=\footnotesize, text=black!80},
  qty/.style   ={font=\footnotesize, text=black!80},% tensor dims, AP numbers
  tsr/.style ={font=\scriptsize\ttfamily, text=black!70},% B x C x H x W (O5)
  tag/.style   ={font=\scriptsize, inner sep=1.6pt, rounded corners=1pt},
  % panel letters (B11): \footnotesize\bfseries at the panel's top-left
  panl/.style  ={font=\footnotesize\bfseries, text=black!85, inner sep=1pt},
  % ---- arrows (O6): one data-flow style, one auxiliary/optional style ----
  flow/.style  ={-{Stealth[length=1.7mm,width=1.35mm]}, draw=flowink,
                 line width=0.65pt},
  flowt/.style ={-{Stealth[length=1.7mm,width=1.35mm]}, draw=flowink,
                 line width=0.65pt, rounded corners=1.6pt},
  aux/.style   ={-{Stealth[length=1.5mm,width=1.2mm]}, draw=black!45,
                 line width=0.5pt, densely dashed},
  auxw/.style  ={-{Stealth[length=1.5mm,width=1.2mm]}, draw=warnred!75,
                 line width=0.5pt, densely dashed},
  rail/.style  ={draw=black!30, line width=0.5pt},
  brace/.style ={decorate, decoration={brace, amplitude=3.2pt, raise=1pt},
                 draw=black!50, line width=0.5pt},
  % ---- ONE leader style for every magnifier cone in the figure set --------
  % v4: the zoom cones used to be drawn in the generic dashed `aux' style, so
  % a sight line was typographically the same object as an ordinary auxiliary
  % arrow.  A cone is now its own object: it carries the hue of the locator
  % rectangle it starts from, no arrowhead (it is a correspondence, not a
  % flow), and it always runs corner-to-corner from the locator box to the
  % MAGNIFIED RASTER, never to the enclosing panel frame.
  zoomray/.style={draw=opsred!55, line width=0.5pt,
                  dash pattern=on 1.3pt off 1.1pt},
  zoombox/.style={draw=opsred, line width=0.7pt},
}

% =========================================================================
% \hbarrow{x0}{y}{track}{frac}{colour}{weight}{label}
%   One horizontal bar on a printed 0..1 track, with the numeral set at the
%   bar end.  v4: fig1's AP progression, figm_tta (c), figm_fusion's result
%   panel and figm_rerank (b) all drew this same object with four slightly
%   different geometries.  They now share one macro, so the reader learns the
%   idiom once and meets it four times -- a small multiple across figures.
%   The empty track is always drawn, so a short bar still shows its scale.
%
%   DEPRECATED FOR METHOD FIGURES (docs/FIGURE_REDESIGN_SPEC.md S5, items 15
%   and 16, rule H1).  \hbarrow / \hbarrowr / \baxis print an accuracy number
%   inside the figure, which the v6 redraw forbids: fig1_framework, figm_bpr
%   and figm_fusion no longer call them and their numbers live in Tables IV
%   and VI.  The macros are KEPT, not deleted, because the parked single-
%   column data panel of Fig. 5 (figm_rerank, permutation importance) is a
%   data figure, where a measured bar length is evidence rather than a
%   result.  Same status, same reason: warn/, badge/ + badgeops/ + badgetrn/,
%   \hull and \mplane -- no v6 method figure uses any of them.
% =========================================================================
\newcommand{\hbarrow}[7]{%
  \draw[black!15, line width=0.5pt] (#1,#2) -- ({#1+#3},#2);
  \draw[#5, line width=#6] (#1,#2) -- ({#1+(#3)*(#4)},#2);
  % the numeral gets a paper backing so the empty part of the track does not
  % run into it and read as a stray rule after the digits
  \node[qty, anchor=west, inner sep=1.6pt, fill=white]
    at ({#1+(#3)*(#4)},#2) {#7};}

% \hbarrowr{x0}{y}{track}{frac}{colour}{weight}{label}
%   As \hbarrow, but the numeral is set in a fixed column just past the end
%   of the track instead of at the bar tip.  Use it where several groups sit
%   side by side on a SHORT track: it aligns the numerals into columns across
%   the groups, and it stops the leftover empty track of a short bar from
%   printing as a stray rule immediately after the digits -- at 1.25 cm of
%   track a 27.7/75 bar leaves 3 mm of rule past the numeral, which reads as
%   an en dash and can be mistaken for a range.
\newcommand{\hbarrowr}[7]{%
  \draw[black!15, line width=0.5pt] (#1,#2) -- ({#1+#3},#2);
  \draw[#5, line width=#6] (#1,#2) -- ({#1+(#3)*(#4)},#2);
  \node[qty, anchor=west, inner sep=1.6pt] at ({#1+(#3)+0.04},#2) {#7};}

% \baxis{x0}{y}{track}{lo}{hi}{unit} : the tick pair + caption for a \hbarrow
%   group, drawn once per group rather than once per bar.
\newcommand{\baxis}[6]{%
  \draw[black!45, line width=0.5pt] (#1,#2) -- (#1,{#2-0.05})
                                    ({#1+#3},#2) -- ({#1+#3},{#2-0.05});
  \node[sub, anchor=north, inner sep=1pt] at (#1,{#2-0.05}) {#4};
  \node[sub, anchor=north, inner sep=1pt] at ({#1+#3},{#2-0.05}) {#5};
  \node[sub, anchor=west, inner sep=1pt] at ({#1+#3+0.14},{#2-0.16}) {#6};}

% =========================================================================
% GEOMETRIC ISOMORPHISM (v3).  Tensor shape is DRAWN, never only written.
%
%  G1 SPATIAL EXTENT -> DRAWN SIZE, one scale for the whole figure set:
%       1 mm of drawn edge per 40 pixels of tensor extent,
%       i.e. drawn side [cm] = side [px] / 400   ->  \pxc{800} = 2.00 cm.
%       800 -> 20.0 mm | 200 -> 5.00 mm | 100 -> 2.50 mm |  50 -> 1.25 mm
%        25 ->  0.63 mm | 640 -> 16.0 mm | 960 -> 24.0 mm |1088 -> 27.2 mm
%        64 ->  1.60 mm | 128 ->  3.20 mm
%       A stride-4 map is therefore exactly 8x the linear size of a
%       stride-32 map, as it must be.
%  G2 CHANNEL COUNT -> SLAB DEPTH (isometric extrusion, 30 degrees):
%       depth [cm] = 0.04 + 0.055 * log2(C/256)
%       256 -> 0.40 mm | 512 -> 0.95 mm | 1024 -> 1.50 mm | 2048 -> 2.05 mm
%       i.e. +0.55 mm of depth per channel doubling.
%  G3 CATEGORY -> GLYPH SHAPE (readable in grayscale, no label needed):
%       spatial feature map (B,C,H,W) ......... extruded slab   \fmslab
%       token set (N x C, no spatial extent) .. column of bars  \tokencol
%       mask stack (N x H x W) ................ offset planes   \maskstack
%       scalar score .......................... short bar       \scorebar
%       feature vector (d cells) .............. strip of cells  (drawn inline)
%  G4 A LABELLED ZOOM INSET may magnify a sub-scene by an explicit factor z;
%       inside it G1/G2 hold multiplied by z.  Every macro takes z as its
%       last argument (z=1 for the global scale).
% =========================================================================
\newcommand{\pxc}[1]{(#1)/400}     % tensor side in px -> drawn side in cm

% =========================================================================
% G5 REAL DATA (v3).  Wherever a glyph stands for a tensor we actually own,
%     the glyph's FRONT FACE carries that tensor's real raster, dumped by
%     carrier/dump_figure_assets.py with provenance in assets/manifest.json.
%     The vector chrome (frames, arrows, labels, extrusion) stays vector, so
%     type remains selectable and the fonts stay embedded -- the bitmap-panel
%     + vector-chrome architecture Nature's artwork guide prescribes.
%     Resolution guard: every embedded raster must clear 300 dpi at its drawn
%     size and we clear 600 dpi everywhere (chip.png 800 px at 20 mm = 1016
%     dpi).  Feature maps are embedded from the nearest-neighbour _xN variants
%     so the PDF renderer cannot smooth pixels the tensor does not have.
%
% \realpanel{x}{y}{sidepx}{zoom}{file}{framespec} : a real raster, framed and
%     sized by G1 exactly as a schematic glyph of the same tensor would be.
% =========================================================================
\newcommand{\realpanel}[6]{%
  \begin{scope}[shift={(#1,#2)}]
    \pgfmathsetmacro{\rpH}{0.5*#4*\pxc{#3}}
    \pgfmathsetmacro{\rpS}{2*\rpH}
    \node[inner sep=0pt, outer sep=0pt] at (0,0)
      {\includegraphics[width=\rpS cm, height=\rpS cm]{\assetdir/#5}};
    \draw[#6] (-\rpH,-\rpH) rectangle (\rpH,\rpH);
  \end{scope}}

% \realpanelh{...} : the same panel HORIZONTALLY MIRRORED.  Used by the TTA
%     figure so the hflip half of each view pair is a real flip of the real
%     raster, not a redrawn mirror of a schematic.
\newcommand{\realpanelh}[6]{%
  \begin{scope}[shift={(#1,#2)}]
    \pgfmathsetmacro{\rpH}{0.5*#4*\pxc{#3}}
    \pgfmathsetmacro{\rpS}{2*\rpH}
    \node[inner sep=0pt, outer sep=0pt] at (0,0)
      {\reflectbox{\includegraphics[width=\rpS cm, height=\rpS cm]{\assetdir/#5}}};
    \draw[#6] (-\rpH,-\rpH) rectangle (\rpH,\rpH);
  \end{scope}}

% --- \fmslab{x}{y}{sidepx}{C}{colour}{zoom} : a spatial feature map -------
\newcommand{\fmslab}[6]{\fmslabr{#1}{#2}{#3}{#4}{#5}{#6}{}}

% --- \fmslabr{...}{file} : the same slab with a REAL raster on the front
%     face.  Pass an empty 7th argument for the flat-tinted variant.  The
%     extruded top/right faces stay flat tinted, so G2 (depth = channel
%     count) and the single light source (top brightest, right middle,
%     front darkest) are preserved.
\newcommand{\fmslabr}[7]{%
  \begin{scope}[shift={(#1,#2)}]
    \pgfmathsetmacro{\sbH}{0.5*#6*\pxc{#3}}
    \pgfmathsetmacro{\sbD}{#6*(0.04+0.055*ln((#4)/256)/ln(2))}
    \pgfmathsetmacro{\sbX}{0.866*\sbD}
    \pgfmathsetmacro{\sbY}{0.500*\sbD}
    \pgfmathsetmacro{\sbW}{2*\sbH}
    \fill[#5!30] (-\sbH,\sbH) -- (-\sbH+\sbX,\sbH+\sbY)
                 -- (\sbH+\sbX,\sbH+\sbY) -- (\sbH,\sbH) -- cycle;
    \fill[#5!20] (\sbH,\sbH) -- (\sbH+\sbX,\sbH+\sbY)
                 -- (\sbH+\sbX,-\sbH+\sbY) -- (\sbH,-\sbH) -- cycle;
    \fill[#5!11] (-\sbH,-\sbH) rectangle (\sbH,\sbH);
    \edef\fmfile{#7}%
    \ifx\fmfile\empty\else
      \node[inner sep=0pt, outer sep=0pt] at (0,0)
        {\includegraphics[width=\sbW cm, height=\sbW cm]{\assetdir/#7}};
    \fi
    \draw[#5!75, line width=0.5pt]
      (-\sbH,-\sbH) rectangle (\sbH,\sbH)
      (-\sbH,\sbH) -- (-\sbH+\sbX,\sbH+\sbY) -- (\sbH+\sbX,\sbH+\sbY) -- (\sbH,\sbH)
      (\sbH+\sbX,\sbH+\sbY) -- (\sbH+\sbX,-\sbH+\sbY) -- (\sbH,-\sbH);
  \end{scope}}

% --- \tokencol{x}{y}{nbars}{width}{colour} : an N x C token set -----------
%     Deliberately NOT a slab: flat bars, visible gaps, no extrusion, so a
%     query set can never be mistaken for a feature map.
\newcommand{\tokencol}[5]{%
  \begin{scope}[shift={(#1,#2)}]
    \pgfmathsetmacro{\tkP}{0.098}
    \pgfmathsetmacro{\tkT}{0.056}
    \pgfmathsetmacro{\tkW}{0.5*#4}
    \pgfmathtruncatemacro{\tkN}{#3-1}
    \pgfmathsetmacro{\tkY}{0.5*(\tkN*\tkP+\tkT)}
    \foreach \tki in {0,...,\tkN} {
      \pgfmathsetmacro{\tkc}{\tkY-\tki*\tkP}
      \fill[#5!13] (-\tkW,\tkc-\tkT) rectangle (\tkW,\tkc);
      \draw[#5!75, line width=0.5pt] (-\tkW,\tkc-\tkT) rectangle (\tkW,\tkc);
    }
  \end{scope}}

% --- \maskstack{x}{y}{sidepx}{k}{colour}{zoom} : k planes of H x W --------
\newcommand{\maskstack}[6]{\maskstackr{#1}{#2}{#3}{#4}{#5}{#6}{}}

% --- \maskstackr{...}{file} : the same pile with the FRONT plane carrying a
%     real raster; the planes behind stay schematic so the "stack of N"
%     reading survives.  Empty 7th argument = the all-schematic variant.
\newcommand{\maskstackr}[7]{%
  \begin{scope}[shift={(#1,#2)}]
    \pgfmathsetmacro{\msH}{0.5*#6*\pxc{#3}}
    \pgfmathsetmacro{\msO}{0.30*\msH/2.5}
    \pgfmathtruncatemacro{\msN}{#4-1}
    \pgfmathsetmacro{\msS}{\msH/0.92}
    \pgfmathsetmacro{\msW}{2*\msH}
    \edef\msfile{#7}%
    \foreach \msi in {\msN,...,0} {
      \pgfmathsetmacro{\msd}{\msi*\msO}
      \pgfmathtruncatemacro{\msreal}{(\msi==0) ? 1 : 0}
      \edef\mshasimg{0}%
      \ifx\msfile\empty\else\ifnum\msreal=1 \edef\mshasimg{1}\fi\fi
      \begin{scope}[shift={(\msd,\msd)}]
        \fill[white] (-\msH,-\msH) rectangle (\msH,\msH);
        \fill[#5!8]  (-\msH,-\msH) rectangle (\msH,\msH);
        \ifnum\mshasimg=1
          \node[inner sep=0pt, outer sep=0pt] at (0,0)
            {\includegraphics[width=\msW cm, height=\msW cm]{\assetdir/#7}};
        \fi
        \draw[#5!60, line width=0.5pt] (-\msH,-\msH) rectangle (\msH,\msH);
        \ifnum\mshasimg=0
          \begin{scope}[scale=\msS]
            \fill[#5!35] \hull;
            \draw[#5!90, line width=0.5pt] \hull;
          \end{scope}
        \fi
      \end{scope}}
  \end{scope}}

% --- \scorebar{x}{y}{h}{colour} : one scalar, as a short upright bar ------
\newcommand{\scorebar}[4]{%
  \begin{scope}[shift={(#1,#2)}]
    \draw[black!35, line width=0.5pt] (-0.10,0) rectangle (0.10,0.62);
    \fill[#4!55] (-0.10,0) rectangle (0.10,#3);
    \draw[#4!90, line width=0.5pt] (-0.10,0) rectangle (0.10,#3);
  \end{scope}}

% =========================================================================
% Reusable glyphs.  All are TikZ pics: \pic[scale=s] at (x,y) {name};
% =========================================================================

% --- lock: marks any component whose weights we never update -------------
\tikzset{lock/.pic={
  \fill[frozgray] (-0.105,-0.105) rectangle (0.105,0.055);
  \draw[frozgray, line width=0.65pt]
    (-0.060,0.055) -- (-0.060,0.105)
    arc[start angle=180, end angle=0, radius=0.060] -- (0.060,0.055);
  \fill[white] (0,-0.022) circle[radius=0.020];
}}

% --- planes: 3 offset stacked planes = a multi-scale feature pyramid -----
\tikzset{planes/.pic={
  \foreach \i/\o in {2/0.20, 1/0.10, 0/0.00} {
    \fill[datablue!\the\numexpr 10+8*\i\relax]
      (\o-0.24,\o-0.30) rectangle (\o+0.24,\o+0.30);
    \draw[datablue!70, line width=0.5pt]
      (\o-0.24,\o-0.30) rectangle (\o+0.24,\o+0.30);
  }
}}

% --- plate: a single feature plane (thin vertical plate, O4) -------------
\tikzset{plate/.pic={
  \fill[datablue!14] (-0.10,-0.34) rectangle (0.10,0.34);
  \draw[datablue!70, line width=0.5pt] (-0.10,-0.34) rectangle (0.10,0.34);
  \draw[datablue!45, line width=0.5pt] (-0.10,0.20) -- (0.10,0.20);
  \draw[datablue!45, line width=0.5pt] (-0.10,-0.20) -- (0.10,-0.20);
}}

% --- ckpt: an isometric slab = one stored checkpoint / weight vector -----
\tikzset{ckpt/.pic={
  \fill[frozgray!18] (-0.30,-0.13) rectangle (0.24,0.13);
  \fill[frozgray!32] (-0.30,0.13) -- (-0.24,0.21) -- (0.30,0.21) -- (0.24,0.13) -- cycle;
  \fill[frozgray!26] (0.24,-0.13) -- (0.30,-0.05) -- (0.30,0.21) -- (0.24,0.13) -- cycle;
  \draw[frozgray, line width=0.5pt]
    (-0.30,-0.13) rectangle (0.24,0.13);
  \draw[frozgray, line width=0.5pt]
    (-0.30,0.13) -- (-0.24,0.21) -- (0.30,0.21) -- (0.24,0.13)
    (0.30,0.21) -- (0.30,-0.05) -- (0.24,-0.13);
}}

% --- masked: pictogram of an output instance mask (hull on light field) --
\tikzset{maskicon/.pic={
  \fill[datablue!5] (-0.30,-0.30) rectangle (0.30,0.30);
  \draw[datablue!60, line width=0.5pt] (-0.30,-0.30) rectangle (0.30,0.30);
  \begin{scope}[rotate=28]
    \fill[datablue!45] (0,0.05) ellipse[x radius=0.235, y radius=0.062];
    \draw[datablue, line width=0.5pt] (0,0.05) ellipse[x radius=0.235, y radius=0.062];
  \end{scope}
  \begin{scope}[rotate=-12]
    \fill[datablue!45] (0.02,-0.19) ellipse[x radius=0.145, y radius=0.045];
    \draw[datablue, line width=0.5pt] (0.02,-0.19) ellipse[x radius=0.145, y radius=0.045];
  \end{scope}
}}

% =========================================================================
% \sarchip was REMOVED in v3.  It drew two smooth ellipses each ~2.6% of the
% frame on a bokeh-like field; the real HRSID inshore median instance covers
% 0.028% of the chip (13.4 px on 800x800), so the glyph was ~90x too large and
% contradicted the small-target-in-clutter argument the paper makes.  Every
% call site now embeds the REAL chip raster instead (assets/chip.png,
% assets/view_{640,800,960,1088}.png; see the realpanel macro above and
% assets/manifest.json).
% =========================================================================

% =========================================================================
% \hull{}  -- schematic ship-mask contour, ~1.3 cm wide, centred at origin.
% Used by the BPR and fusion figures so every ship outline in the paper is
% literally the same curve.
% =========================================================================
\newcommand{\hull}{plot[smooth cycle, tension=0.72] coordinates
  {(0.72,0.00) (0.34,0.26) (-0.22,0.33) (-0.64,0.17) (-0.62,-0.12)
   (-0.10,-0.30) (0.42,-0.22)}}
 in the preamble of a standalone figure.
%
% -------------------------------------------------------------------------
% DESIGN NOTES — conventions extracted from an accepted TGRS manuscript
% (Gao et al., "BBox-free ... SAR ship instance segmentation", TGRS 2024;
%  local copy: scratchpad/gao_bboxfree_tgrs24.pdf).  Observations taken from
%  its Fig. 2 (overall architecture, p. 5), Fig. 5 (FD module, p. 8) and
%  Fig. 6 (conceptual mask panel, p. 8):
%
%  O1  REGION PARTITION.  The architecture figure is carved into 3-4 named
%      regions by DOTTED rectangles with no fill; the region name is set in
%      bold OUTSIDE the rectangle ("Feature Extractor", "Classification
%      Branch", "Localization Branch", "Center Point Extractor").  Nothing
%      is nested more than two levels deep.
%  O2  COLOR CARRIES MEANING, NOT DECORATION.  Only ~5 hues appear, all as
%      pale fills with black text: light gray = ordinary/standard module,
%      salmon-orange = the module the paper contributes (Feature Decoupling
%      Module, Center Point Extractor), yellow = a distinguished operation
%      (Feature Fusion), light blue / light pink = tensors.  Saturated color
%      is reserved for the few glyphs that must be found first.
%  O3  LEGEND PANEL.  Both Fig. 2 and Fig. 5 carry an explicit legend inside
%      the figure (dotted box top-left in Fig. 2, dashed box right in Fig. 5)
%      that decodes arrow colors and single-letter glyph nodes.  The body of
%      the diagram is kept uncluttered by pushing vocabulary into the legend.
%  O4  TENSOR ICONOGRAPHY.  Feature maps are drawn as thin vertical plates,
%      multi-scale features as 3 offset stacked planes, channel vectors as
%      segmented 3-D bars, the backbone as an isometric slab.  Real SAR
%      chips / heat maps / GT overlays are pasted directly into the flow.
%  O5  TENSOR NAMES SIT BESIDE THE GLYPH, NEVER INSIDE IT, in math italic
%      (F_out, F_mask, F_param, F_H_zip ...), placed above-right of the plate.
%  O6  ARROW DISCIPLINE.  Data flow is a thin gray/black line with a small
%      solid triangular head; routing is orthogonal (Manhattan) with a shared
%      trunk when one source fans out to several targets, so lines never
%      cross.  Semantically special transitions get a THICK COLORED BLOCK
%      arrow instead (2x upsample / 2x downsample / connect), i.e. arrow
%      style is itself a semantic channel.  Long skip paths are routed as
%      rails around the outside of the diagram (Fig. 5 top and bottom).
%  O7  TYPOGRAPHY.  One size for block labels, one smaller size for
%      sub-labels, math only in math mode; no stretched or rotated text
%      except a rare 90-degree axis label.
%  O8  ASPECT RATIO.  Double-column method figures are wide and short:
%      Fig. 2 is 530 x 200 pt (2.65:1), Fig. 5 is 530 x 215 pt (2.5:1).
%      They are laid on a grid with equal gutters between parallel branches.
%  O9  CAPTION DIVISION OF LABOR.  The caption names the parts and decodes
%      abbreviations/colors; it does not re-derive the mechanism, which is
%      left to the body text.  The figure is readable without the caption.
%  O10 CONCEPTUAL PANELS (Fig. 6) use dashed sub-panels with a centered
%      label underneath, a pale blue blob for the true object and a pale red
%      overlay for the error region -- schematic, not photorealistic.
%
% -------------------------------------------------------------------------
% WHAT WE ADOPT / WHERE WE DEVIATE
%  + adopt O1 (dotted named regions), O3 (in-figure legend), O4/O5 (tensor
%    glyphs with the symbol outside), O6 (one flow style + Manhattan
%    routing), O8 (wide-short double-column ratios), O10 (schematic panels).
%  ~ deviate on O2: instead of hue-as-identity we bind hue to ROLE in the
%    contribution -- frozen / training-free operator / trained corrector /
%    data -- because the entire argument of this article is which parts are
%    touched and which are not.
%  ~ deviate on the sans-serif block labels of Fig. 2: we set every figure in
%    the body Times face (as the same paper's Fig. 6 does), so figure text
%    and caption text share one face.
%  + every figure is authored at its EXACT printed width (18.2 cm double
%    column = 516 pt, 8.89 cm single column = 252 pt) and included at
%    scale 1, so \footnotesize is a true 8 pt and \scriptsize a true 7 pt on
%    the page in ALL figures.  Use \figcanvasd / \figcanvasc below.
%
% -------------------------------------------------------------------------
% SEMANTIC PALETTE (Okabe-Ito colorblind-safe base; redundant encoding by
% border weight, dash pattern and badge glyph so it survives grayscale):
%   froz  gray   + lock glyph        frozen, not modified by us  (Stage A)
%   ops   vermil + solid heavy rule  our TRAINING-FREE operator  (rung 1)
%   trn   green  + dashed rule       our OPTIONAL TRAINED corrector (2-3)
%   data  blue   + thin rule         data, tensors, I/O
%   warn  red    + dotted rule, small  failure mode / negative control
% =========================================================================

% Real rasters are embedded on the front faces of the tensor glyphs (see
% \realpanel / \fmslab below).  graphicx must be loaded in the preamble, which
% is where this file is \input.
\usepackage{graphicx}
% Single switch for the raster payload: \def\assetdir{assets_light}
% before \begin{document} swaps every embedded raster for its
% print-sized ~600 dpi copy (see figs/make_light_assets.py).  The
% default is the native dump, 965-1016 dpi at the drawn size.
\providecommand{\assetdir}{assets}
\usetikzlibrary{positioning,arrows.meta,fit,backgrounds,calc,shapes.geometric,
                decorations.pathreplacing,decorations.markings,patterns}

% ---------------- palette ----------------
\definecolor{opsred}{HTML}{D55E00}   % Okabe-Ito vermillion  -- training-free ops
\definecolor{trngreen}{HTML}{009E73} % Okabe-Ito bluish green-- trained correctors
\definecolor{datablue}{HTML}{0072B2} % Okabe-Ito blue        -- data / tensors
\definecolor{frozgray}{HTML}{6E6E6E} % neutral               -- frozen base
\definecolor{flowink}{HTML}{262626}  % near black            -- data-flow arrows
\definecolor{warnred}{HTML}{9E2B2B}  % muted brick           -- failure modes
\definecolor{sarink}{HTML}{2B2B33}   % SAR chip background

% ---------------- exact-width canvases ----------------
% \figcanvasd{H} : double column, 18.2 cm = 516 pt = \textwidth  (scale 1.0)
% \figcanvasc{H} : single column,  8.89 cm = 252 pt = \columnwidth
\newcommand{\figcanvasd}[1]{\useasboundingbox (0,0) rectangle (18.2,#1);}
\newcommand{\figcanvasc}[1]{\useasboundingbox (0,0) rectangle (8.89,#1);}

\tikzset{
  % ---- picture-level defaults: ONE font size system for all five figures ----
  fig/.style={font=\footnotesize, line width=0.65pt,
              >={Stealth[length=1.7mm,width=1.35mm]},
              every node/.append style={align=center, inner sep=2.4pt}},
  % ---- blocks: role -> color + redundant border encoding ----
  % Fill tints were raised in v3: at 4-8% the five roles were separated by an
  % RGB distance of only 6.6-21 (out of 441) and all landed on gray 243-248,
  % i.e. the fill channel was consuming ink and returning nothing.  The border
  % weight / dash pattern / badge remain the PRIMARY, grayscale-safe channel.
  ops/.style   ={draw=opsred,   fill=opsred!16,   line width=0.9pt,
                 rounded corners=1.6pt, inner sep=3pt, align=center},
  trn/.style   ={draw=trngreen, fill=trngreen!15, line width=0.9pt,
                 dash pattern=on 1.7pt off 1.1pt,
                 rounded corners=1.6pt, inner sep=3pt, align=center},
  froz/.style  ={draw=frozgray, fill=black!8,     line width=0.7pt,
                 rounded corners=1.6pt, inner sep=3pt, align=center},
  data/.style  ={draw=datablue, fill=datablue!14, line width=0.55pt,
                 rounded corners=1.6pt, inner sep=3pt, align=center},
  plain/.style ={draw=black!45, fill=white,       line width=0.55pt,
                 rounded corners=1.6pt, inner sep=3pt, align=center},
  % v4: radius 1.4 -> 1.6 pt, the one block style that did not share the
  % system's corner radius; weight 0.5 -> 0.6 pt so the failure chips are not
  % the faintest object on the page while being the ones a reader must find.
  % The dotted dash pattern stays as the grayscale-safe channel.
  warn/.style  ={draw=warnred,  fill=warnred!9,   line width=0.6pt,
                 dash pattern=on 1.1pt off 1.1pt,
                 rounded corners=1.6pt, inner sep=3pt, align=center},
  gate/.style  ={draw=opsred,   fill=opsred!16,   line width=0.9pt,
                 diamond, aspect=2.15, inner sep=1pt, align=center},
  % ---- named region (O1): dotted, unfilled, label set outside ----
  grp/.style   ={draw=black!45, densely dotted, line width=0.7pt,
                 rounded corners=2pt, inner sep=0pt},
  grptitle/.style={font=\footnotesize\bfseries, text=black!80, inner sep=1.5pt},
  % ---- step badges (redundant, grayscale-safe stage identity) ----
  badge/.style ={circle, inner sep=0pt, minimum size=3.3mm, text=white,
                 font=\scriptsize\bfseries},
  badgeops/.style={badge, fill=opsred},
  badgetrn/.style={badge, fill=trngreen},
  % ---- typography (O7): exactly three text sizes in the whole figure set ----
  % v3: 7 pt at 55% black was the worst corner of the design space (minimum
  % size AND minimum contrast) -- sub-labels are now 7 pt at 70% black, and
  % anything carrying a QUANTITY is promoted to 8 pt via \qty / \subk.
  ttl/.style   ={font=\footnotesize\bfseries},
  lbl/.style   ={font=\footnotesize},
  sub/.style   ={font=\scriptsize, text=black!70},
  subk/.style  ={font=\footnotesize, text=black!80},
  qty/.style   ={font=\footnotesize, text=black!80},% tensor dims, AP numbers
  tsr/.style ={font=\scriptsize\ttfamily, text=black!70},% B x C x H x W (O5)
  tag/.style   ={font=\scriptsize, inner sep=1.6pt, rounded corners=1pt},
  % panel letters (B11): \footnotesize\bfseries at the panel's top-left
  panl/.style  ={font=\footnotesize\bfseries, text=black!85, inner sep=1pt},
  % ---- arrows (O6): one data-flow style, one auxiliary/optional style ----
  flow/.style  ={-{Stealth[length=1.7mm,width=1.35mm]}, draw=flowink,
                 line width=0.65pt},
  flowt/.style ={-{Stealth[length=1.7mm,width=1.35mm]}, draw=flowink,
                 line width=0.65pt, rounded corners=1.6pt},
  aux/.style   ={-{Stealth[length=1.5mm,width=1.2mm]}, draw=black!45,
                 line width=0.5pt, densely dashed},
  auxw/.style  ={-{Stealth[length=1.5mm,width=1.2mm]}, draw=warnred!75,
                 line width=0.5pt, densely dashed},
  rail/.style  ={draw=black!30, line width=0.5pt},
  brace/.style ={decorate, decoration={brace, amplitude=3.2pt, raise=1pt},
                 draw=black!50, line width=0.5pt},
  % ---- ONE leader style for every magnifier cone in the figure set --------
  % v4: the zoom cones used to be drawn in the generic dashed `aux' style, so
  % a sight line was typographically the same object as an ordinary auxiliary
  % arrow.  A cone is now its own object: it carries the hue of the locator
  % rectangle it starts from, no arrowhead (it is a correspondence, not a
  % flow), and it always runs corner-to-corner from the locator box to the
  % MAGNIFIED RASTER, never to the enclosing panel frame.
  zoomray/.style={draw=opsred!55, line width=0.5pt,
                  dash pattern=on 1.3pt off 1.1pt},
  zoombox/.style={draw=opsred, line width=0.7pt},
}

% =========================================================================
% \hbarrow{x0}{y}{track}{frac}{colour}{weight}{label}
%   One horizontal bar on a printed 0..1 track, with the numeral set at the
%   bar end.  v4: fig1's AP progression, figm_tta (c), figm_fusion's result
%   panel and figm_rerank (b) all drew this same object with four slightly
%   different geometries.  They now share one macro, so the reader learns the
%   idiom once and meets it four times -- a small multiple across figures.
%   The empty track is always drawn, so a short bar still shows its scale.
%
%   DEPRECATED FOR METHOD FIGURES (docs/FIGURE_REDESIGN_SPEC.md S5, items 15
%   and 16, rule H1).  \hbarrow / \hbarrowr / \baxis print an accuracy number
%   inside the figure, which the v6 redraw forbids: fig1_framework, figm_bpr
%   and figm_fusion no longer call them and their numbers live in Tables IV
%   and VI.  The macros are KEPT, not deleted, because the parked single-
%   column data panel of Fig. 5 (figm_rerank, permutation importance) is a
%   data figure, where a measured bar length is evidence rather than a
%   result.  Same status, same reason: warn/, badge/ + badgeops/ + badgetrn/,
%   \hull and \mplane -- no v6 method figure uses any of them.
% =========================================================================
\newcommand{\hbarrow}[7]{%
  \draw[black!15, line width=0.5pt] (#1,#2) -- ({#1+#3},#2);
  \draw[#5, line width=#6] (#1,#2) -- ({#1+(#3)*(#4)},#2);
  % the numeral gets a paper backing so the empty part of the track does not
  % run into it and read as a stray rule after the digits
  \node[qty, anchor=west, inner sep=1.6pt, fill=white]
    at ({#1+(#3)*(#4)},#2) {#7};}

% \hbarrowr{x0}{y}{track}{frac}{colour}{weight}{label}
%   As \hbarrow, but the numeral is set in a fixed column just past the end
%   of the track instead of at the bar tip.  Use it where several groups sit
%   side by side on a SHORT track: it aligns the numerals into columns across
%   the groups, and it stops the leftover empty track of a short bar from
%   printing as a stray rule immediately after the digits -- at 1.25 cm of
%   track a 27.7/75 bar leaves 3 mm of rule past the numeral, which reads as
%   an en dash and can be mistaken for a range.
\newcommand{\hbarrowr}[7]{%
  \draw[black!15, line width=0.5pt] (#1,#2) -- ({#1+#3},#2);
  \draw[#5, line width=#6] (#1,#2) -- ({#1+(#3)*(#4)},#2);
  \node[qty, anchor=west, inner sep=1.6pt] at ({#1+(#3)+0.04},#2) {#7};}

% \baxis{x0}{y}{track}{lo}{hi}{unit} : the tick pair + caption for a \hbarrow
%   group, drawn once per group rather than once per bar.
\newcommand{\baxis}[6]{%
  \draw[black!45, line width=0.5pt] (#1,#2) -- (#1,{#2-0.05})
                                    ({#1+#3},#2) -- ({#1+#3},{#2-0.05});
  \node[sub, anchor=north, inner sep=1pt] at (#1,{#2-0.05}) {#4};
  \node[sub, anchor=north, inner sep=1pt] at ({#1+#3},{#2-0.05}) {#5};
  \node[sub, anchor=west, inner sep=1pt] at ({#1+#3+0.14},{#2-0.16}) {#6};}

% =========================================================================
% GEOMETRIC ISOMORPHISM (v3).  Tensor shape is DRAWN, never only written.
%
%  G1 SPATIAL EXTENT -> DRAWN SIZE, one scale for the whole figure set:
%       1 mm of drawn edge per 40 pixels of tensor extent,
%       i.e. drawn side [cm] = side [px] / 400   ->  \pxc{800} = 2.00 cm.
%       800 -> 20.0 mm | 200 -> 5.00 mm | 100 -> 2.50 mm |  50 -> 1.25 mm
%        25 ->  0.63 mm | 640 -> 16.0 mm | 960 -> 24.0 mm |1088 -> 27.2 mm
%        64 ->  1.60 mm | 128 ->  3.20 mm
%       A stride-4 map is therefore exactly 8x the linear size of a
%       stride-32 map, as it must be.
%  G2 CHANNEL COUNT -> SLAB DEPTH (isometric extrusion, 30 degrees):
%       depth [cm] = 0.04 + 0.055 * log2(C/256)
%       256 -> 0.40 mm | 512 -> 0.95 mm | 1024 -> 1.50 mm | 2048 -> 2.05 mm
%       i.e. +0.55 mm of depth per channel doubling.
%  G3 CATEGORY -> GLYPH SHAPE (readable in grayscale, no label needed):
%       spatial feature map (B,C,H,W) ......... extruded slab   \fmslab
%       token set (N x C, no spatial extent) .. column of bars  \tokencol
%       mask stack (N x H x W) ................ offset planes   \maskstack
%       scalar score .......................... short bar       \scorebar
%       feature vector (d cells) .............. strip of cells  (drawn inline)
%  G4 A LABELLED ZOOM INSET may magnify a sub-scene by an explicit factor z;
%       inside it G1/G2 hold multiplied by z.  Every macro takes z as its
%       last argument (z=1 for the global scale).
% =========================================================================
\newcommand{\pxc}[1]{(#1)/400}     % tensor side in px -> drawn side in cm

% =========================================================================
% G5 REAL DATA (v3).  Wherever a glyph stands for a tensor we actually own,
%     the glyph's FRONT FACE carries that tensor's real raster, dumped by
%     carrier/dump_figure_assets.py with provenance in assets/manifest.json.
%     The vector chrome (frames, arrows, labels, extrusion) stays vector, so
%     type remains selectable and the fonts stay embedded -- the bitmap-panel
%     + vector-chrome architecture Nature's artwork guide prescribes.
%     Resolution guard: every embedded raster must clear 300 dpi at its drawn
%     size and we clear 600 dpi everywhere (chip.png 800 px at 20 mm = 1016
%     dpi).  Feature maps are embedded from the nearest-neighbour _xN variants
%     so the PDF renderer cannot smooth pixels the tensor does not have.
%
% \realpanel{x}{y}{sidepx}{zoom}{file}{framespec} : a real raster, framed and
%     sized by G1 exactly as a schematic glyph of the same tensor would be.
% =========================================================================
\newcommand{\realpanel}[6]{%
  \begin{scope}[shift={(#1,#2)}]
    \pgfmathsetmacro{\rpH}{0.5*#4*\pxc{#3}}
    \pgfmathsetmacro{\rpS}{2*\rpH}
    \node[inner sep=0pt, outer sep=0pt] at (0,0)
      {\includegraphics[width=\rpS cm, height=\rpS cm]{\assetdir/#5}};
    \draw[#6] (-\rpH,-\rpH) rectangle (\rpH,\rpH);
  \end{scope}}

% \realpanelh{...} : the same panel HORIZONTALLY MIRRORED.  Used by the TTA
%     figure so the hflip half of each view pair is a real flip of the real
%     raster, not a redrawn mirror of a schematic.
\newcommand{\realpanelh}[6]{%
  \begin{scope}[shift={(#1,#2)}]
    \pgfmathsetmacro{\rpH}{0.5*#4*\pxc{#3}}
    \pgfmathsetmacro{\rpS}{2*\rpH}
    \node[inner sep=0pt, outer sep=0pt] at (0,0)
      {\reflectbox{\includegraphics[width=\rpS cm, height=\rpS cm]{\assetdir/#5}}};
    \draw[#6] (-\rpH,-\rpH) rectangle (\rpH,\rpH);
  \end{scope}}

% --- \fmslab{x}{y}{sidepx}{C}{colour}{zoom} : a spatial feature map -------
\newcommand{\fmslab}[6]{\fmslabr{#1}{#2}{#3}{#4}{#5}{#6}{}}

% --- \fmslabr{...}{file} : the same slab with a REAL raster on the front
%     face.  Pass an empty 7th argument for the flat-tinted variant.  The
%     extruded top/right faces stay flat tinted, so G2 (depth = channel
%     count) and the single light source (top brightest, right middle,
%     front darkest) are preserved.
\newcommand{\fmslabr}[7]{%
  \begin{scope}[shift={(#1,#2)}]
    \pgfmathsetmacro{\sbH}{0.5*#6*\pxc{#3}}
    \pgfmathsetmacro{\sbD}{#6*(0.04+0.055*ln((#4)/256)/ln(2))}
    \pgfmathsetmacro{\sbX}{0.866*\sbD}
    \pgfmathsetmacro{\sbY}{0.500*\sbD}
    \pgfmathsetmacro{\sbW}{2*\sbH}
    \fill[#5!30] (-\sbH,\sbH) -- (-\sbH+\sbX,\sbH+\sbY)
                 -- (\sbH+\sbX,\sbH+\sbY) -- (\sbH,\sbH) -- cycle;
    \fill[#5!20] (\sbH,\sbH) -- (\sbH+\sbX,\sbH+\sbY)
                 -- (\sbH+\sbX,-\sbH+\sbY) -- (\sbH,-\sbH) -- cycle;
    \fill[#5!11] (-\sbH,-\sbH) rectangle (\sbH,\sbH);
    \edef\fmfile{#7}%
    \ifx\fmfile\empty\else
      \node[inner sep=0pt, outer sep=0pt] at (0,0)
        {\includegraphics[width=\sbW cm, height=\sbW cm]{\assetdir/#7}};
    \fi
    \draw[#5!75, line width=0.5pt]
      (-\sbH,-\sbH) rectangle (\sbH,\sbH)
      (-\sbH,\sbH) -- (-\sbH+\sbX,\sbH+\sbY) -- (\sbH+\sbX,\sbH+\sbY) -- (\sbH,\sbH)
      (\sbH+\sbX,\sbH+\sbY) -- (\sbH+\sbX,-\sbH+\sbY) -- (\sbH,-\sbH);
  \end{scope}}

% --- \tokencol{x}{y}{nbars}{width}{colour} : an N x C token set -----------
%     Deliberately NOT a slab: flat bars, visible gaps, no extrusion, so a
%     query set can never be mistaken for a feature map.
\newcommand{\tokencol}[5]{%
  \begin{scope}[shift={(#1,#2)}]
    \pgfmathsetmacro{\tkP}{0.098}
    \pgfmathsetmacro{\tkT}{0.056}
    \pgfmathsetmacro{\tkW}{0.5*#4}
    \pgfmathtruncatemacro{\tkN}{#3-1}
    \pgfmathsetmacro{\tkY}{0.5*(\tkN*\tkP+\tkT)}
    \foreach \tki in {0,...,\tkN} {
      \pgfmathsetmacro{\tkc}{\tkY-\tki*\tkP}
      \fill[#5!13] (-\tkW,\tkc-\tkT) rectangle (\tkW,\tkc);
      \draw[#5!75, line width=0.5pt] (-\tkW,\tkc-\tkT) rectangle (\tkW,\tkc);
    }
  \end{scope}}

% --- \maskstack{x}{y}{sidepx}{k}{colour}{zoom} : k planes of H x W --------
\newcommand{\maskstack}[6]{\maskstackr{#1}{#2}{#3}{#4}{#5}{#6}{}}

% --- \maskstackr{...}{file} : the same pile with the FRONT plane carrying a
%     real raster; the planes behind stay schematic so the "stack of N"
%     reading survives.  Empty 7th argument = the all-schematic variant.
\newcommand{\maskstackr}[7]{%
  \begin{scope}[shift={(#1,#2)}]
    \pgfmathsetmacro{\msH}{0.5*#6*\pxc{#3}}
    \pgfmathsetmacro{\msO}{0.30*\msH/2.5}
    \pgfmathtruncatemacro{\msN}{#4-1}
    \pgfmathsetmacro{\msS}{\msH/0.92}
    \pgfmathsetmacro{\msW}{2*\msH}
    \edef\msfile{#7}%
    \foreach \msi in {\msN,...,0} {
      \pgfmathsetmacro{\msd}{\msi*\msO}
      \pgfmathtruncatemacro{\msreal}{(\msi==0) ? 1 : 0}
      \edef\mshasimg{0}%
      \ifx\msfile\empty\else\ifnum\msreal=1 \edef\mshasimg{1}\fi\fi
      \begin{scope}[shift={(\msd,\msd)}]
        \fill[white] (-\msH,-\msH) rectangle (\msH,\msH);
        \fill[#5!8]  (-\msH,-\msH) rectangle (\msH,\msH);
        \ifnum\mshasimg=1
          \node[inner sep=0pt, outer sep=0pt] at (0,0)
            {\includegraphics[width=\msW cm, height=\msW cm]{\assetdir/#7}};
        \fi
        \draw[#5!60, line width=0.5pt] (-\msH,-\msH) rectangle (\msH,\msH);
        \ifnum\mshasimg=0
          \begin{scope}[scale=\msS]
            \fill[#5!35] \hull;
            \draw[#5!90, line width=0.5pt] \hull;
          \end{scope}
        \fi
      \end{scope}}
  \end{scope}}

% --- \scorebar{x}{y}{h}{colour} : one scalar, as a short upright bar ------
\newcommand{\scorebar}[4]{%
  \begin{scope}[shift={(#1,#2)}]
    \draw[black!35, line width=0.5pt] (-0.10,0) rectangle (0.10,0.62);
    \fill[#4!55] (-0.10,0) rectangle (0.10,#3);
    \draw[#4!90, line width=0.5pt] (-0.10,0) rectangle (0.10,#3);
  \end{scope}}

% =========================================================================
% Reusable glyphs.  All are TikZ pics: \pic[scale=s] at (x,y) {name};
% =========================================================================

% --- lock: marks any component whose weights we never update -------------
\tikzset{lock/.pic={
  \fill[frozgray] (-0.105,-0.105) rectangle (0.105,0.055);
  \draw[frozgray, line width=0.65pt]
    (-0.060,0.055) -- (-0.060,0.105)
    arc[start angle=180, end angle=0, radius=0.060] -- (0.060,0.055);
  \fill[white] (0,-0.022) circle[radius=0.020];
}}

% --- planes: 3 offset stacked planes = a multi-scale feature pyramid -----
\tikzset{planes/.pic={
  \foreach \i/\o in {2/0.20, 1/0.10, 0/0.00} {
    \fill[datablue!\the\numexpr 10+8*\i\relax]
      (\o-0.24,\o-0.30) rectangle (\o+0.24,\o+0.30);
    \draw[datablue!70, line width=0.5pt]
      (\o-0.24,\o-0.30) rectangle (\o+0.24,\o+0.30);
  }
}}

% --- plate: a single feature plane (thin vertical plate, O4) -------------
\tikzset{plate/.pic={
  \fill[datablue!14] (-0.10,-0.34) rectangle (0.10,0.34);
  \draw[datablue!70, line width=0.5pt] (-0.10,-0.34) rectangle (0.10,0.34);
  \draw[datablue!45, line width=0.5pt] (-0.10,0.20) -- (0.10,0.20);
  \draw[datablue!45, line width=0.5pt] (-0.10,-0.20) -- (0.10,-0.20);
}}

% --- ckpt: an isometric slab = one stored checkpoint / weight vector -----
\tikzset{ckpt/.pic={
  \fill[frozgray!18] (-0.30,-0.13) rectangle (0.24,0.13);
  \fill[frozgray!32] (-0.30,0.13) -- (-0.24,0.21) -- (0.30,0.21) -- (0.24,0.13) -- cycle;
  \fill[frozgray!26] (0.24,-0.13) -- (0.30,-0.05) -- (0.30,0.21) -- (0.24,0.13) -- cycle;
  \draw[frozgray, line width=0.5pt]
    (-0.30,-0.13) rectangle (0.24,0.13);
  \draw[frozgray, line width=0.5pt]
    (-0.30,0.13) -- (-0.24,0.21) -- (0.30,0.21) -- (0.24,0.13)
    (0.30,0.21) -- (0.30,-0.05) -- (0.24,-0.13);
}}

% --- masked: pictogram of an output instance mask (hull on light field) --
\tikzset{maskicon/.pic={
  \fill[datablue!5] (-0.30,-0.30) rectangle (0.30,0.30);
  \draw[datablue!60, line width=0.5pt] (-0.30,-0.30) rectangle (0.30,0.30);
  \begin{scope}[rotate=28]
    \fill[datablue!45] (0,0.05) ellipse[x radius=0.235, y radius=0.062];
    \draw[datablue, line width=0.5pt] (0,0.05) ellipse[x radius=0.235, y radius=0.062];
  \end{scope}
  \begin{scope}[rotate=-12]
    \fill[datablue!45] (0.02,-0.19) ellipse[x radius=0.145, y radius=0.045];
    \draw[datablue, line width=0.5pt] (0.02,-0.19) ellipse[x radius=0.145, y radius=0.045];
  \end{scope}
}}

% =========================================================================
% \sarchip was REMOVED in v3.  It drew two smooth ellipses each ~2.6% of the
% frame on a bokeh-like field; the real HRSID inshore median instance covers
% 0.028% of the chip (13.4 px on 800x800), so the glyph was ~90x too large and
% contradicted the small-target-in-clutter argument the paper makes.  Every
% call site now embeds the REAL chip raster instead (assets/chip.png,
% assets/view_{640,800,960,1088}.png; see the realpanel macro above and
% assets/manifest.json).
% =========================================================================

% =========================================================================
% \hull{}  -- schematic ship-mask contour, ~1.3 cm wide, centred at origin.
% Used by the BPR and fusion figures so every ship outline in the paper is
% literally the same curve.
% =========================================================================
\newcommand{\hull}{plot[smooth cycle, tension=0.72] coordinates
  {(0.72,0.00) (0.34,0.26) (-0.22,0.33) (-0.64,0.17) (-0.62,-0.12)
   (-0.10,-0.30) (0.42,-0.22)}}
 in the preamble of a standalone figure.
%
% -------------------------------------------------------------------------
% DESIGN NOTES — conventions extracted from an accepted TGRS manuscript
% (Gao et al., "BBox-free ... SAR ship instance segmentation", TGRS 2024;
%  local copy: scratchpad/gao_bboxfree_tgrs24.pdf).  Observations taken from
%  its Fig. 2 (overall architecture, p. 5), Fig. 5 (FD module, p. 8) and
%  Fig. 6 (conceptual mask panel, p. 8):
%
%  O1  REGION PARTITION.  The architecture figure is carved into 3-4 named
%      regions by DOTTED rectangles with no fill; the region name is set in
%      bold OUTSIDE the rectangle ("Feature Extractor", "Classification
%      Branch", "Localization Branch", "Center Point Extractor").  Nothing
%      is nested more than two levels deep.
%  O2  COLOR CARRIES MEANING, NOT DECORATION.  Only ~5 hues appear, all as
%      pale fills with black text: light gray = ordinary/standard module,
%      salmon-orange = the module the paper contributes (Feature Decoupling
%      Module, Center Point Extractor), yellow = a distinguished operation
%      (Feature Fusion), light blue / light pink = tensors.  Saturated color
%      is reserved for the few glyphs that must be found first.
%  O3  LEGEND PANEL.  Both Fig. 2 and Fig. 5 carry an explicit legend inside
%      the figure (dotted box top-left in Fig. 2, dashed box right in Fig. 5)
%      that decodes arrow colors and single-letter glyph nodes.  The body of
%      the diagram is kept uncluttered by pushing vocabulary into the legend.
%  O4  TENSOR ICONOGRAPHY.  Feature maps are drawn as thin vertical plates,
%      multi-scale features as 3 offset stacked planes, channel vectors as
%      segmented 3-D bars, the backbone as an isometric slab.  Real SAR
%      chips / heat maps / GT overlays are pasted directly into the flow.
%  O5  TENSOR NAMES SIT BESIDE THE GLYPH, NEVER INSIDE IT, in math italic
%      (F_out, F_mask, F_param, F_H_zip ...), placed above-right of the plate.
%  O6  ARROW DISCIPLINE.  Data flow is a thin gray/black line with a small
%      solid triangular head; routing is orthogonal (Manhattan) with a shared
%      trunk when one source fans out to several targets, so lines never
%      cross.  Semantically special transitions get a THICK COLORED BLOCK
%      arrow instead (2x upsample / 2x downsample / connect), i.e. arrow
%      style is itself a semantic channel.  Long skip paths are routed as
%      rails around the outside of the diagram (Fig. 5 top and bottom).
%  O7  TYPOGRAPHY.  One size for block labels, one smaller size for
%      sub-labels, math only in math mode; no stretched or rotated text
%      except a rare 90-degree axis label.
%  O8  ASPECT RATIO.  Double-column method figures are wide and short:
%      Fig. 2 is 530 x 200 pt (2.65:1), Fig. 5 is 530 x 215 pt (2.5:1).
%      They are laid on a grid with equal gutters between parallel branches.
%  O9  CAPTION DIVISION OF LABOR.  The caption names the parts and decodes
%      abbreviations/colors; it does not re-derive the mechanism, which is
%      left to the body text.  The figure is readable without the caption.
%  O10 CONCEPTUAL PANELS (Fig. 6) use dashed sub-panels with a centered
%      label underneath, a pale blue blob for the true object and a pale red
%      overlay for the error region -- schematic, not photorealistic.
%
% -------------------------------------------------------------------------
% WHAT WE ADOPT / WHERE WE DEVIATE
%  + adopt O1 (dotted named regions), O3 (in-figure legend), O4/O5 (tensor
%    glyphs with the symbol outside), O6 (one flow style + Manhattan
%    routing), O8 (wide-short double-column ratios), O10 (schematic panels).
%  ~ deviate on O2: instead of hue-as-identity we bind hue to ROLE in the
%    contribution -- frozen / training-free operator / trained corrector /
%    data -- because the entire argument of this article is which parts are
%    touched and which are not.
%  ~ deviate on the sans-serif block labels of Fig. 2: we set every figure in
%    the body Times face (as the same paper's Fig. 6 does), so figure text
%    and caption text share one face.
%  + every figure is authored at its EXACT printed width (18.2 cm double
%    column = 516 pt, 8.89 cm single column = 252 pt) and included at
%    scale 1, so \footnotesize is a true 8 pt and \scriptsize a true 7 pt on
%    the page in ALL figures.  Use \figcanvasd / \figcanvasc below.
%
% -------------------------------------------------------------------------
% SEMANTIC PALETTE (Okabe-Ito colorblind-safe base; redundant encoding by
% border weight, dash pattern and badge glyph so it survives grayscale):
%   froz  gray   + lock glyph        frozen, not modified by us  (Stage A)
%   ops   vermil + solid heavy rule  our TRAINING-FREE operator  (rung 1)
%   trn   green  + dashed rule       our OPTIONAL TRAINED corrector (2-3)
%   data  blue   + thin rule         data, tensors, I/O
%   warn  red    + dotted rule, small  failure mode / negative control
% =========================================================================

% Real rasters are embedded on the front faces of the tensor glyphs (see
% \realpanel / \fmslab below).  graphicx must be loaded in the preamble, which
% is where this file is \input.
\usepackage{graphicx}
% Single switch for the raster payload: \def\assetdir{assets_light}
% before \begin{document} swaps every embedded raster for its
% print-sized ~600 dpi copy (see figs/make_light_assets.py).  The
% default is the native dump, 965-1016 dpi at the drawn size.
\providecommand{\assetdir}{assets}
\usetikzlibrary{positioning,arrows.meta,fit,backgrounds,calc,shapes.geometric,
                decorations.pathreplacing,decorations.markings,patterns}

% ---------------- palette ----------------
\definecolor{opsred}{HTML}{D55E00}   % Okabe-Ito vermillion  -- training-free ops
\definecolor{trngreen}{HTML}{009E73} % Okabe-Ito bluish green-- trained correctors
\definecolor{datablue}{HTML}{0072B2} % Okabe-Ito blue        -- data / tensors
\definecolor{frozgray}{HTML}{6E6E6E} % neutral               -- frozen base
\definecolor{flowink}{HTML}{262626}  % near black            -- data-flow arrows
\definecolor{warnred}{HTML}{9E2B2B}  % muted brick           -- failure modes
\definecolor{sarink}{HTML}{2B2B33}   % SAR chip background

% ---------------- exact-width canvases ----------------
% \figcanvasd{H} : double column, 18.2 cm = 516 pt = \textwidth  (scale 1.0)
% \figcanvasc{H} : single column,  8.89 cm = 252 pt = \columnwidth
\newcommand{\figcanvasd}[1]{\useasboundingbox (0,0) rectangle (18.2,#1);}
\newcommand{\figcanvasc}[1]{\useasboundingbox (0,0) rectangle (8.89,#1);}

\tikzset{
  % ---- picture-level defaults: ONE font size system for all five figures ----
  fig/.style={font=\footnotesize, line width=0.65pt,
              >={Stealth[length=1.7mm,width=1.35mm]},
              every node/.append style={align=center, inner sep=2.4pt}},
  % ---- blocks: role -> color + redundant border encoding ----
  % Fill tints were raised in v3: at 4-8% the five roles were separated by an
  % RGB distance of only 6.6-21 (out of 441) and all landed on gray 243-248,
  % i.e. the fill channel was consuming ink and returning nothing.  The border
  % weight / dash pattern / badge remain the PRIMARY, grayscale-safe channel.
  ops/.style   ={draw=opsred,   fill=opsred!16,   line width=0.9pt,
                 rounded corners=1.6pt, inner sep=3pt, align=center},
  trn/.style   ={draw=trngreen, fill=trngreen!15, line width=0.9pt,
                 dash pattern=on 1.7pt off 1.1pt,
                 rounded corners=1.6pt, inner sep=3pt, align=center},
  froz/.style  ={draw=frozgray, fill=black!8,     line width=0.7pt,
                 rounded corners=1.6pt, inner sep=3pt, align=center},
  data/.style  ={draw=datablue, fill=datablue!14, line width=0.55pt,
                 rounded corners=1.6pt, inner sep=3pt, align=center},
  plain/.style ={draw=black!45, fill=white,       line width=0.55pt,
                 rounded corners=1.6pt, inner sep=3pt, align=center},
  % v4: radius 1.4 -> 1.6 pt, the one block style that did not share the
  % system's corner radius; weight 0.5 -> 0.6 pt so the failure chips are not
  % the faintest object on the page while being the ones a reader must find.
  % The dotted dash pattern stays as the grayscale-safe channel.
  warn/.style  ={draw=warnred,  fill=warnred!9,   line width=0.6pt,
                 dash pattern=on 1.1pt off 1.1pt,
                 rounded corners=1.6pt, inner sep=3pt, align=center},
  gate/.style  ={draw=opsred,   fill=opsred!16,   line width=0.9pt,
                 diamond, aspect=2.15, inner sep=1pt, align=center},
  % ---- named region (O1): dotted, unfilled, label set outside ----
  grp/.style   ={draw=black!45, densely dotted, line width=0.7pt,
                 rounded corners=2pt, inner sep=0pt},
  grptitle/.style={font=\footnotesize\bfseries, text=black!80, inner sep=1.5pt},
  % ---- step badges (redundant, grayscale-safe stage identity) ----
  badge/.style ={circle, inner sep=0pt, minimum size=3.3mm, text=white,
                 font=\scriptsize\bfseries},
  badgeops/.style={badge, fill=opsred},
  badgetrn/.style={badge, fill=trngreen},
  % ---- typography (O7): exactly three text sizes in the whole figure set ----
  % v3: 7 pt at 55% black was the worst corner of the design space (minimum
  % size AND minimum contrast) -- sub-labels are now 7 pt at 70% black, and
  % anything carrying a QUANTITY is promoted to 8 pt via \qty / \subk.
  ttl/.style   ={font=\footnotesize\bfseries},
  lbl/.style   ={font=\footnotesize},
  sub/.style   ={font=\scriptsize, text=black!70},
  subk/.style  ={font=\footnotesize, text=black!80},
  qty/.style   ={font=\footnotesize, text=black!80},% tensor dims, AP numbers
  tsr/.style ={font=\scriptsize\ttfamily, text=black!70},% B x C x H x W (O5)
  tag/.style   ={font=\scriptsize, inner sep=1.6pt, rounded corners=1pt},
  % panel letters (B11): \footnotesize\bfseries at the panel's top-left
  panl/.style  ={font=\footnotesize\bfseries, text=black!85, inner sep=1pt},
  % ---- arrows (O6): one data-flow style, one auxiliary/optional style ----
  flow/.style  ={-{Stealth[length=1.7mm,width=1.35mm]}, draw=flowink,
                 line width=0.65pt},
  flowt/.style ={-{Stealth[length=1.7mm,width=1.35mm]}, draw=flowink,
                 line width=0.65pt, rounded corners=1.6pt},
  aux/.style   ={-{Stealth[length=1.5mm,width=1.2mm]}, draw=black!45,
                 line width=0.5pt, densely dashed},
  auxw/.style  ={-{Stealth[length=1.5mm,width=1.2mm]}, draw=warnred!75,
                 line width=0.5pt, densely dashed},
  rail/.style  ={draw=black!30, line width=0.5pt},
  brace/.style ={decorate, decoration={brace, amplitude=3.2pt, raise=1pt},
                 draw=black!50, line width=0.5pt},
  % ---- ONE leader style for every magnifier cone in the figure set --------
  % v4: the zoom cones used to be drawn in the generic dashed `aux' style, so
  % a sight line was typographically the same object as an ordinary auxiliary
  % arrow.  A cone is now its own object: it carries the hue of the locator
  % rectangle it starts from, no arrowhead (it is a correspondence, not a
  % flow), and it always runs corner-to-corner from the locator box to the
  % MAGNIFIED RASTER, never to the enclosing panel frame.
  zoomray/.style={draw=opsred!55, line width=0.5pt,
                  dash pattern=on 1.3pt off 1.1pt},
  zoombox/.style={draw=opsred, line width=0.7pt},
}

% =========================================================================
% \hbarrow{x0}{y}{track}{frac}{colour}{weight}{label}
%   One horizontal bar on a printed 0..1 track, with the numeral set at the
%   bar end.  v4: fig1's AP progression, figm_tta (c), figm_fusion's result
%   panel and figm_rerank (b) all drew this same object with four slightly
%   different geometries.  They now share one macro, so the reader learns the
%   idiom once and meets it four times -- a small multiple across figures.
%   The empty track is always drawn, so a short bar still shows its scale.
%
%   DEPRECATED FOR METHOD FIGURES (docs/FIGURE_REDESIGN_SPEC.md S5, items 15
%   and 16, rule H1).  \hbarrow / \hbarrowr / \baxis print an accuracy number
%   inside the figure, which the v6 redraw forbids: fig1_framework, figm_bpr
%   and figm_fusion no longer call them and their numbers live in Tables IV
%   and VI.  The macros are KEPT, not deleted, because the parked single-
%   column data panel of Fig. 5 (figm_rerank, permutation importance) is a
%   data figure, where a measured bar length is evidence rather than a
%   result.  Same status, same reason: warn/, badge/ + badgeops/ + badgetrn/,
%   \hull and \mplane -- no v6 method figure uses any of them.
% =========================================================================
\newcommand{\hbarrow}[7]{%
  \draw[black!15, line width=0.5pt] (#1,#2) -- ({#1+#3},#2);
  \draw[#5, line width=#6] (#1,#2) -- ({#1+(#3)*(#4)},#2);
  % the numeral gets a paper backing so the empty part of the track does not
  % run into it and read as a stray rule after the digits
  \node[qty, anchor=west, inner sep=1.6pt, fill=white]
    at ({#1+(#3)*(#4)},#2) {#7};}

% \hbarrowr{x0}{y}{track}{frac}{colour}{weight}{label}
%   As \hbarrow, but the numeral is set in a fixed column just past the end
%   of the track instead of at the bar tip.  Use it where several groups sit
%   side by side on a SHORT track: it aligns the numerals into columns across
%   the groups, and it stops the leftover empty track of a short bar from
%   printing as a stray rule immediately after the digits -- at 1.25 cm of
%   track a 27.7/75 bar leaves 3 mm of rule past the numeral, which reads as
%   an en dash and can be mistaken for a range.
\newcommand{\hbarrowr}[7]{%
  \draw[black!15, line width=0.5pt] (#1,#2) -- ({#1+#3},#2);
  \draw[#5, line width=#6] (#1,#2) -- ({#1+(#3)*(#4)},#2);
  \node[qty, anchor=west, inner sep=1.6pt] at ({#1+(#3)+0.04},#2) {#7};}

% \baxis{x0}{y}{track}{lo}{hi}{unit} : the tick pair + caption for a \hbarrow
%   group, drawn once per group rather than once per bar.
\newcommand{\baxis}[6]{%
  \draw[black!45, line width=0.5pt] (#1,#2) -- (#1,{#2-0.05})
                                    ({#1+#3},#2) -- ({#1+#3},{#2-0.05});
  \node[sub, anchor=north, inner sep=1pt] at (#1,{#2-0.05}) {#4};
  \node[sub, anchor=north, inner sep=1pt] at ({#1+#3},{#2-0.05}) {#5};
  \node[sub, anchor=west, inner sep=1pt] at ({#1+#3+0.14},{#2-0.16}) {#6};}

% =========================================================================
% GEOMETRIC ISOMORPHISM (v3).  Tensor shape is DRAWN, never only written.
%
%  G1 SPATIAL EXTENT -> DRAWN SIZE, one scale for the whole figure set:
%       1 mm of drawn edge per 40 pixels of tensor extent,
%       i.e. drawn side [cm] = side [px] / 400   ->  \pxc{800} = 2.00 cm.
%       800 -> 20.0 mm | 200 -> 5.00 mm | 100 -> 2.50 mm |  50 -> 1.25 mm
%        25 ->  0.63 mm | 640 -> 16.0 mm | 960 -> 24.0 mm |1088 -> 27.2 mm
%        64 ->  1.60 mm | 128 ->  3.20 mm
%       A stride-4 map is therefore exactly 8x the linear size of a
%       stride-32 map, as it must be.
%  G2 CHANNEL COUNT -> SLAB DEPTH (isometric extrusion, 30 degrees):
%       depth [cm] = 0.04 + 0.055 * log2(C/256)
%       256 -> 0.40 mm | 512 -> 0.95 mm | 1024 -> 1.50 mm | 2048 -> 2.05 mm
%       i.e. +0.55 mm of depth per channel doubling.
%  G3 CATEGORY -> GLYPH SHAPE (readable in grayscale, no label needed):
%       spatial feature map (B,C,H,W) ......... extruded slab   \fmslab
%       token set (N x C, no spatial extent) .. column of bars  \tokencol
%       mask stack (N x H x W) ................ offset planes   \maskstack
%       scalar score .......................... short bar       \scorebar
%       feature vector (d cells) .............. strip of cells  (drawn inline)
%  G4 A LABELLED ZOOM INSET may magnify a sub-scene by an explicit factor z;
%       inside it G1/G2 hold multiplied by z.  Every macro takes z as its
%       last argument (z=1 for the global scale).
% =========================================================================
\newcommand{\pxc}[1]{(#1)/400}     % tensor side in px -> drawn side in cm

% =========================================================================
% G5 REAL DATA (v3).  Wherever a glyph stands for a tensor we actually own,
%     the glyph's FRONT FACE carries that tensor's real raster, dumped by
%     carrier/dump_figure_assets.py with provenance in assets/manifest.json.
%     The vector chrome (frames, arrows, labels, extrusion) stays vector, so
%     type remains selectable and the fonts stay embedded -- the bitmap-panel
%     + vector-chrome architecture Nature's artwork guide prescribes.
%     Resolution guard: every embedded raster must clear 300 dpi at its drawn
%     size and we clear 600 dpi everywhere (chip.png 800 px at 20 mm = 1016
%     dpi).  Feature maps are embedded from the nearest-neighbour _xN variants
%     so the PDF renderer cannot smooth pixels the tensor does not have.
%
% \realpanel{x}{y}{sidepx}{zoom}{file}{framespec} : a real raster, framed and
%     sized by G1 exactly as a schematic glyph of the same tensor would be.
% =========================================================================
\newcommand{\realpanel}[6]{%
  \begin{scope}[shift={(#1,#2)}]
    \pgfmathsetmacro{\rpH}{0.5*#4*\pxc{#3}}
    \pgfmathsetmacro{\rpS}{2*\rpH}
    \node[inner sep=0pt, outer sep=0pt] at (0,0)
      {\includegraphics[width=\rpS cm, height=\rpS cm]{\assetdir/#5}};
    \draw[#6] (-\rpH,-\rpH) rectangle (\rpH,\rpH);
  \end{scope}}

% \realpanelh{...} : the same panel HORIZONTALLY MIRRORED.  Used by the TTA
%     figure so the hflip half of each view pair is a real flip of the real
%     raster, not a redrawn mirror of a schematic.
\newcommand{\realpanelh}[6]{%
  \begin{scope}[shift={(#1,#2)}]
    \pgfmathsetmacro{\rpH}{0.5*#4*\pxc{#3}}
    \pgfmathsetmacro{\rpS}{2*\rpH}
    \node[inner sep=0pt, outer sep=0pt] at (0,0)
      {\reflectbox{\includegraphics[width=\rpS cm, height=\rpS cm]{\assetdir/#5}}};
    \draw[#6] (-\rpH,-\rpH) rectangle (\rpH,\rpH);
  \end{scope}}

% --- \fmslab{x}{y}{sidepx}{C}{colour}{zoom} : a spatial feature map -------
\newcommand{\fmslab}[6]{\fmslabr{#1}{#2}{#3}{#4}{#5}{#6}{}}

% --- \fmslabr{...}{file} : the same slab with a REAL raster on the front
%     face.  Pass an empty 7th argument for the flat-tinted variant.  The
%     extruded top/right faces stay flat tinted, so G2 (depth = channel
%     count) and the single light source (top brightest, right middle,
%     front darkest) are preserved.
\newcommand{\fmslabr}[7]{%
  \begin{scope}[shift={(#1,#2)}]
    \pgfmathsetmacro{\sbH}{0.5*#6*\pxc{#3}}
    \pgfmathsetmacro{\sbD}{#6*(0.04+0.055*ln((#4)/256)/ln(2))}
    \pgfmathsetmacro{\sbX}{0.866*\sbD}
    \pgfmathsetmacro{\sbY}{0.500*\sbD}
    \pgfmathsetmacro{\sbW}{2*\sbH}
    \fill[#5!30] (-\sbH,\sbH) -- (-\sbH+\sbX,\sbH+\sbY)
                 -- (\sbH+\sbX,\sbH+\sbY) -- (\sbH,\sbH) -- cycle;
    \fill[#5!20] (\sbH,\sbH) -- (\sbH+\sbX,\sbH+\sbY)
                 -- (\sbH+\sbX,-\sbH+\sbY) -- (\sbH,-\sbH) -- cycle;
    \fill[#5!11] (-\sbH,-\sbH) rectangle (\sbH,\sbH);
    \edef\fmfile{#7}%
    \ifx\fmfile\empty\else
      \node[inner sep=0pt, outer sep=0pt] at (0,0)
        {\includegraphics[width=\sbW cm, height=\sbW cm]{\assetdir/#7}};
    \fi
    \draw[#5!75, line width=0.5pt]
      (-\sbH,-\sbH) rectangle (\sbH,\sbH)
      (-\sbH,\sbH) -- (-\sbH+\sbX,\sbH+\sbY) -- (\sbH+\sbX,\sbH+\sbY) -- (\sbH,\sbH)
      (\sbH+\sbX,\sbH+\sbY) -- (\sbH+\sbX,-\sbH+\sbY) -- (\sbH,-\sbH);
  \end{scope}}

% --- \tokencol{x}{y}{nbars}{width}{colour} : an N x C token set -----------
%     Deliberately NOT a slab: flat bars, visible gaps, no extrusion, so a
%     query set can never be mistaken for a feature map.
\newcommand{\tokencol}[5]{%
  \begin{scope}[shift={(#1,#2)}]
    \pgfmathsetmacro{\tkP}{0.098}
    \pgfmathsetmacro{\tkT}{0.056}
    \pgfmathsetmacro{\tkW}{0.5*#4}
    \pgfmathtruncatemacro{\tkN}{#3-1}
    \pgfmathsetmacro{\tkY}{0.5*(\tkN*\tkP+\tkT)}
    \foreach \tki in {0,...,\tkN} {
      \pgfmathsetmacro{\tkc}{\tkY-\tki*\tkP}
      \fill[#5!13] (-\tkW,\tkc-\tkT) rectangle (\tkW,\tkc);
      \draw[#5!75, line width=0.5pt] (-\tkW,\tkc-\tkT) rectangle (\tkW,\tkc);
    }
  \end{scope}}

% --- \maskstack{x}{y}{sidepx}{k}{colour}{zoom} : k planes of H x W --------
\newcommand{\maskstack}[6]{\maskstackr{#1}{#2}{#3}{#4}{#5}{#6}{}}

% --- \maskstackr{...}{file} : the same pile with the FRONT plane carrying a
%     real raster; the planes behind stay schematic so the "stack of N"
%     reading survives.  Empty 7th argument = the all-schematic variant.
\newcommand{\maskstackr}[7]{%
  \begin{scope}[shift={(#1,#2)}]
    \pgfmathsetmacro{\msH}{0.5*#6*\pxc{#3}}
    \pgfmathsetmacro{\msO}{0.30*\msH/2.5}
    \pgfmathtruncatemacro{\msN}{#4-1}
    \pgfmathsetmacro{\msS}{\msH/0.92}
    \pgfmathsetmacro{\msW}{2*\msH}
    \edef\msfile{#7}%
    \foreach \msi in {\msN,...,0} {
      \pgfmathsetmacro{\msd}{\msi*\msO}
      \pgfmathtruncatemacro{\msreal}{(\msi==0) ? 1 : 0}
      \edef\mshasimg{0}%
      \ifx\msfile\empty\else\ifnum\msreal=1 \edef\mshasimg{1}\fi\fi
      \begin{scope}[shift={(\msd,\msd)}]
        \fill[white] (-\msH,-\msH) rectangle (\msH,\msH);
        \fill[#5!8]  (-\msH,-\msH) rectangle (\msH,\msH);
        \ifnum\mshasimg=1
          \node[inner sep=0pt, outer sep=0pt] at (0,0)
            {\includegraphics[width=\msW cm, height=\msW cm]{\assetdir/#7}};
        \fi
        \draw[#5!60, line width=0.5pt] (-\msH,-\msH) rectangle (\msH,\msH);
        \ifnum\mshasimg=0
          \begin{scope}[scale=\msS]
            \fill[#5!35] \hull;
            \draw[#5!90, line width=0.5pt] \hull;
          \end{scope}
        \fi
      \end{scope}}
  \end{scope}}

% --- \scorebar{x}{y}{h}{colour} : one scalar, as a short upright bar ------
\newcommand{\scorebar}[4]{%
  \begin{scope}[shift={(#1,#2)}]
    \draw[black!35, line width=0.5pt] (-0.10,0) rectangle (0.10,0.62);
    \fill[#4!55] (-0.10,0) rectangle (0.10,#3);
    \draw[#4!90, line width=0.5pt] (-0.10,0) rectangle (0.10,#3);
  \end{scope}}

% =========================================================================
% Reusable glyphs.  All are TikZ pics: \pic[scale=s] at (x,y) {name};
% =========================================================================

% --- lock: marks any component whose weights we never update -------------
\tikzset{lock/.pic={
  \fill[frozgray] (-0.105,-0.105) rectangle (0.105,0.055);
  \draw[frozgray, line width=0.65pt]
    (-0.060,0.055) -- (-0.060,0.105)
    arc[start angle=180, end angle=0, radius=0.060] -- (0.060,0.055);
  \fill[white] (0,-0.022) circle[radius=0.020];
}}

% --- planes: 3 offset stacked planes = a multi-scale feature pyramid -----
\tikzset{planes/.pic={
  \foreach \i/\o in {2/0.20, 1/0.10, 0/0.00} {
    \fill[datablue!\the\numexpr 10+8*\i\relax]
      (\o-0.24,\o-0.30) rectangle (\o+0.24,\o+0.30);
    \draw[datablue!70, line width=0.5pt]
      (\o-0.24,\o-0.30) rectangle (\o+0.24,\o+0.30);
  }
}}

% --- plate: a single feature plane (thin vertical plate, O4) -------------
\tikzset{plate/.pic={
  \fill[datablue!14] (-0.10,-0.34) rectangle (0.10,0.34);
  \draw[datablue!70, line width=0.5pt] (-0.10,-0.34) rectangle (0.10,0.34);
  \draw[datablue!45, line width=0.5pt] (-0.10,0.20) -- (0.10,0.20);
  \draw[datablue!45, line width=0.5pt] (-0.10,-0.20) -- (0.10,-0.20);
}}

% --- ckpt: an isometric slab = one stored checkpoint / weight vector -----
\tikzset{ckpt/.pic={
  \fill[frozgray!18] (-0.30,-0.13) rectangle (0.24,0.13);
  \fill[frozgray!32] (-0.30,0.13) -- (-0.24,0.21) -- (0.30,0.21) -- (0.24,0.13) -- cycle;
  \fill[frozgray!26] (0.24,-0.13) -- (0.30,-0.05) -- (0.30,0.21) -- (0.24,0.13) -- cycle;
  \draw[frozgray, line width=0.5pt]
    (-0.30,-0.13) rectangle (0.24,0.13);
  \draw[frozgray, line width=0.5pt]
    (-0.30,0.13) -- (-0.24,0.21) -- (0.30,0.21) -- (0.24,0.13)
    (0.30,0.21) -- (0.30,-0.05) -- (0.24,-0.13);
}}

% --- masked: pictogram of an output instance mask (hull on light field) --
\tikzset{maskicon/.pic={
  \fill[datablue!5] (-0.30,-0.30) rectangle (0.30,0.30);
  \draw[datablue!60, line width=0.5pt] (-0.30,-0.30) rectangle (0.30,0.30);
  \begin{scope}[rotate=28]
    \fill[datablue!45] (0,0.05) ellipse[x radius=0.235, y radius=0.062];
    \draw[datablue, line width=0.5pt] (0,0.05) ellipse[x radius=0.235, y radius=0.062];
  \end{scope}
  \begin{scope}[rotate=-12]
    \fill[datablue!45] (0.02,-0.19) ellipse[x radius=0.145, y radius=0.045];
    \draw[datablue, line width=0.5pt] (0.02,-0.19) ellipse[x radius=0.145, y radius=0.045];
  \end{scope}
}}

% =========================================================================
% \sarchip was REMOVED in v3.  It drew two smooth ellipses each ~2.6% of the
% frame on a bokeh-like field; the real HRSID inshore median instance covers
% 0.028% of the chip (13.4 px on 800x800), so the glyph was ~90x too large and
% contradicted the small-target-in-clutter argument the paper makes.  Every
% call site now embeds the REAL chip raster instead (assets/chip.png,
% assets/view_{640,800,960,1088}.png; see the realpanel macro above and
% assets/manifest.json).
% =========================================================================

% =========================================================================
% \hull{}  -- schematic ship-mask contour, ~1.3 cm wide, centred at origin.
% Used by the BPR and fusion figures so every ship outline in the paper is
% literally the same curve.
% =========================================================================
\newcommand{\hull}{plot[smooth cycle, tension=0.72] coordinates
  {(0.72,0.00) (0.34,0.26) (-0.22,0.33) (-0.64,0.17) (-0.62,-0.12)
   (-0.10,-0.30) (0.42,-0.22)}}
 in the preamble of a standalone figure.
%
% -------------------------------------------------------------------------
% DESIGN NOTES — conventions extracted from an accepted TGRS manuscript
% (Gao et al., "BBox-free ... SAR ship instance segmentation", TGRS 2024;
%  local copy: scratchpad/gao_bboxfree_tgrs24.pdf).  Observations taken from
%  its Fig. 2 (overall architecture, p. 5), Fig. 5 (FD module, p. 8) and
%  Fig. 6 (conceptual mask panel, p. 8):
%
%  O1  REGION PARTITION.  The architecture figure is carved into 3-4 named
%      regions by DOTTED rectangles with no fill; the region name is set in
%      bold OUTSIDE the rectangle ("Feature Extractor", "Classification
%      Branch", "Localization Branch", "Center Point Extractor").  Nothing
%      is nested more than two levels deep.
%  O2  COLOR CARRIES MEANING, NOT DECORATION.  Only ~5 hues appear, all as
%      pale fills with black text: light gray = ordinary/standard module,
%      salmon-orange = the module the paper contributes (Feature Decoupling
%      Module, Center Point Extractor), yellow = a distinguished operation
%      (Feature Fusion), light blue / light pink = tensors.  Saturated color
%      is reserved for the few glyphs that must be found first.
%  O3  LEGEND PANEL.  Both Fig. 2 and Fig. 5 carry an explicit legend inside
%      the figure (dotted box top-left in Fig. 2, dashed box right in Fig. 5)
%      that decodes arrow colors and single-letter glyph nodes.  The body of
%      the diagram is kept uncluttered by pushing vocabulary into the legend.
%  O4  TENSOR ICONOGRAPHY.  Feature maps are drawn as thin vertical plates,
%      multi-scale features as 3 offset stacked planes, channel vectors as
%      segmented 3-D bars, the backbone as an isometric slab.  Real SAR
%      chips / heat maps / GT overlays are pasted directly into the flow.
%  O5  TENSOR NAMES SIT BESIDE THE GLYPH, NEVER INSIDE IT, in math italic
%      (F_out, F_mask, F_param, F_H_zip ...), placed above-right of the plate.
%  O6  ARROW DISCIPLINE.  Data flow is a thin gray/black line with a small
%      solid triangular head; routing is orthogonal (Manhattan) with a shared
%      trunk when one source fans out to several targets, so lines never
%      cross.  Semantically special transitions get a THICK COLORED BLOCK
%      arrow instead (2x upsample / 2x downsample / connect), i.e. arrow
%      style is itself a semantic channel.  Long skip paths are routed as
%      rails around the outside of the diagram (Fig. 5 top and bottom).
%  O7  TYPOGRAPHY.  One size for block labels, one smaller size for
%      sub-labels, math only in math mode; no stretched or rotated text
%      except a rare 90-degree axis label.
%  O8  ASPECT RATIO.  Double-column method figures are wide and short:
%      Fig. 2 is 530 x 200 pt (2.65:1), Fig. 5 is 530 x 215 pt (2.5:1).
%      They are laid on a grid with equal gutters between parallel branches.
%  O9  CAPTION DIVISION OF LABOR.  The caption names the parts and decodes
%      abbreviations/colors; it does not re-derive the mechanism, which is
%      left to the body text.  The figure is readable without the caption.
%  O10 CONCEPTUAL PANELS (Fig. 6) use dashed sub-panels with a centered
%      label underneath, a pale blue blob for the true object and a pale red
%      overlay for the error region -- schematic, not photorealistic.
%
% -------------------------------------------------------------------------
% WHAT WE ADOPT / WHERE WE DEVIATE
%  + adopt O1 (dotted named regions), O3 (in-figure legend), O4/O5 (tensor
%    glyphs with the symbol outside), O6 (one flow style + Manhattan
%    routing), O8 (wide-short double-column ratios), O10 (schematic panels).
%  ~ deviate on O2: instead of hue-as-identity we bind hue to ROLE in the
%    contribution -- frozen / training-free operator / trained corrector /
%    data -- because the entire argument of this article is which parts are
%    touched and which are not.
%  ~ deviate on the sans-serif block labels of Fig. 2: we set every figure in
%    the body Times face (as the same paper's Fig. 6 does), so figure text
%    and caption text share one face.
%  + every figure is authored at its EXACT printed width (18.2 cm double
%    column = 516 pt, 8.89 cm single column = 252 pt) and included at
%    scale 1, so \footnotesize is a true 8 pt and \scriptsize a true 7 pt on
%    the page in ALL figures.  Use \figcanvasd / \figcanvasc below.
%
% -------------------------------------------------------------------------
% SEMANTIC PALETTE (Okabe-Ito colorblind-safe base; redundant encoding by
% border weight, dash pattern and badge glyph so it survives grayscale):
%   froz  gray   + lock glyph        frozen, not modified by us  (Stage A)
%   ops   vermil + solid heavy rule  our TRAINING-FREE operator  (rung 1)
%   trn   green  + dashed rule       our OPTIONAL TRAINED corrector (2-3)
%   data  blue   + thin rule         data, tensors, I/O
%   warn  red    + dotted rule, small  failure mode / negative control
% =========================================================================

\usetikzlibrary{positioning,arrows.meta,fit,backgrounds,calc,shapes.geometric,
                decorations.pathreplacing,decorations.markings,patterns,
                shadows.blur}

% ---------------- palette ----------------
\definecolor{opsred}{HTML}{D55E00}   % Okabe-Ito vermillion  -- training-free ops
\definecolor{trngreen}{HTML}{009E73} % Okabe-Ito bluish green-- trained correctors
\definecolor{datablue}{HTML}{0072B2} % Okabe-Ito blue        -- data / tensors
\definecolor{frozgray}{HTML}{6E6E6E} % neutral               -- frozen base
\definecolor{flowink}{HTML}{262626}  % near black            -- data-flow arrows
\definecolor{warnred}{HTML}{9E2B2B}  % muted brick           -- failure modes
\definecolor{sarink}{HTML}{2B2B33}   % SAR chip background

% ---------------- exact-width canvases ----------------
% \figcanvasd{H} : double column, 18.2 cm = 516 pt = \textwidth  (scale 1.0)
% \figcanvasc{H} : single column,  8.89 cm = 252 pt = \columnwidth
\newcommand{\figcanvasd}[1]{\useasboundingbox (0,0) rectangle (18.2,#1);}
\newcommand{\figcanvasc}[1]{\useasboundingbox (0,0) rectangle (8.89,#1);}

\tikzset{
  % ---- picture-level defaults: ONE font size system for all five figures ----
  fig/.style={font=\footnotesize, line width=0.65pt,
              >={Stealth[length=1.7mm,width=1.35mm]},
              every node/.append style={align=center, inner sep=2.4pt}},
  % ---- blocks: role -> color + redundant border encoding ----
  ops/.style   ={draw=opsred,   fill=opsred!8,    line width=0.9pt,
                 rounded corners=1.6pt, inner sep=3pt, align=center},
  trn/.style   ={draw=trngreen, fill=trngreen!8,  line width=0.9pt,
                 dash pattern=on 1.7pt off 1.1pt,
                 rounded corners=1.6pt, inner sep=3pt, align=center},
  froz/.style  ={draw=frozgray, fill=black!4,     line width=0.7pt,
                 rounded corners=1.6pt, inner sep=3pt, align=center},
  data/.style  ={draw=datablue, fill=datablue!7,  line width=0.55pt,
                 rounded corners=1.6pt, inner sep=3pt, align=center},
  plain/.style ={draw=black!45, fill=white,       line width=0.55pt,
                 rounded corners=1.6pt, inner sep=3pt, align=center},
  warn/.style  ={draw=warnred,  fill=warnred!4,   line width=0.5pt,
                 dash pattern=on 1.1pt off 1.1pt,
                 rounded corners=1.4pt, inner sep=3pt, align=center},
  gate/.style  ={draw=opsred,   fill=opsred!8,    line width=0.9pt,
                 diamond, aspect=2.15, inner sep=1pt, align=center},
  % ---- named region (O1): dotted, unfilled, label set outside ----
  grp/.style   ={draw=black!45, densely dotted, line width=0.7pt,
                 rounded corners=2pt, inner sep=0pt},
  grptitle/.style={font=\footnotesize\bfseries, text=black!80, inner sep=1.5pt},
  % ---- step badges (redundant, grayscale-safe stage identity) ----
  badge/.style ={circle, inner sep=0pt, minimum size=3.3mm, text=white,
                 font=\scriptsize\bfseries},
  badgeops/.style={badge, fill=opsred},
  badgetrn/.style={badge, fill=trngreen},
  % ---- typography (O7): exactly three text sizes in the whole figure set ----
  ttl/.style   ={font=\footnotesize\bfseries},
  lbl/.style   ={font=\footnotesize},
  sub/.style   ={font=\scriptsize, text=black!55},
  subk/.style  ={font=\scriptsize, text=black!80},
  tsr/.style ={font=\scriptsize\ttfamily, text=black!55},% B x C x H x W (O5)
  tag/.style   ={font=\scriptsize, inner sep=1.6pt, rounded corners=1pt},
  % ---- arrows (O6): one data-flow style, one auxiliary/optional style ----
  flow/.style  ={-{Stealth[length=1.7mm,width=1.35mm]}, draw=flowink,
                 line width=0.65pt},
  flowt/.style ={-{Stealth[length=1.7mm,width=1.35mm]}, draw=flowink,
                 line width=0.65pt, rounded corners=1.6pt},
  aux/.style   ={-{Stealth[length=1.5mm,width=1.2mm]}, draw=black!45,
                 line width=0.5pt, densely dashed},
  auxw/.style  ={-{Stealth[length=1.5mm,width=1.2mm]}, draw=warnred!75,
                 line width=0.5pt, densely dashed},
  rail/.style  ={draw=black!30, line width=0.5pt},
  brace/.style ={decorate, decoration={brace, amplitude=3.2pt, raise=1pt},
                 draw=black!50, line width=0.5pt},
}

% =========================================================================
% GEOMETRIC ISOMORPHISM (v3).  Tensor shape is DRAWN, never only written.
%
%  G1 SPATIAL EXTENT -> DRAWN SIZE, one scale for the whole figure set:
%       1 mm of drawn edge per 40 pixels of tensor extent,
%       i.e. drawn side [cm] = side [px] / 400   ->  \pxc{800} = 2.00 cm.
%       800 -> 20.0 mm | 200 -> 5.00 mm | 100 -> 2.50 mm |  50 -> 1.25 mm
%        25 ->  0.63 mm | 640 -> 16.0 mm | 960 -> 24.0 mm |1088 -> 27.2 mm
%        64 ->  1.60 mm | 128 ->  3.20 mm
%       A stride-4 map is therefore exactly 8x the linear size of a
%       stride-32 map, as it must be.
%  G2 CHANNEL COUNT -> SLAB DEPTH (isometric extrusion, 30 degrees):
%       depth [cm] = 0.04 + 0.055 * log2(C/256)
%       256 -> 0.40 mm | 512 -> 0.95 mm | 1024 -> 1.50 mm | 2048 -> 2.05 mm
%       i.e. +0.55 mm of depth per channel doubling.
%  G3 CATEGORY -> GLYPH SHAPE (readable in grayscale, no label needed):
%       spatial feature map (B,C,H,W) ......... extruded slab   \fmslab
%       token set (N x C, no spatial extent) .. column of bars  \tokencol
%       mask stack (N x H x W) ................ offset planes   \maskstack
%       scalar score .......................... short bar       \scorebar
%       feature vector (d cells) .............. strip of cells  (drawn inline)
%  G4 A LABELLED ZOOM INSET may magnify a sub-scene by an explicit factor z;
%       inside it G1/G2 hold multiplied by z.  Every macro takes z as its
%       last argument (z=1 for the global scale).
% =========================================================================
\newcommand{\pxc}[1]{(#1)/400}     % tensor side in px -> drawn side in cm

% --- \fmslab{x}{y}{sidepx}{C}{colour}{zoom} : a spatial feature map -------
\newcommand{\fmslab}[6]{%
  \begin{scope}[shift={(#1,#2)}]
    \pgfmathsetmacro{\sbH}{0.5*#6*\pxc{#3}}
    \pgfmathsetmacro{\sbD}{#6*(0.04+0.055*ln((#4)/256)/ln(2))}
    \pgfmathsetmacro{\sbX}{0.866*\sbD}
    \pgfmathsetmacro{\sbY}{0.500*\sbD}
    \fill[#5!30] (-\sbH,\sbH) -- (-\sbH+\sbX,\sbH+\sbY)
                 -- (\sbH+\sbX,\sbH+\sbY) -- (\sbH,\sbH) -- cycle;
    \fill[#5!20] (\sbH,\sbH) -- (\sbH+\sbX,\sbH+\sbY)
                 -- (\sbH+\sbX,-\sbH+\sbY) -- (\sbH,-\sbH) -- cycle;
    \fill[#5!11] (-\sbH,-\sbH) rectangle (\sbH,\sbH);
    \draw[#5!75, line width=0.4pt]
      (-\sbH,-\sbH) rectangle (\sbH,\sbH)
      (-\sbH,\sbH) -- (-\sbH+\sbX,\sbH+\sbY) -- (\sbH+\sbX,\sbH+\sbY) -- (\sbH,\sbH)
      (\sbH+\sbX,\sbH+\sbY) -- (\sbH+\sbX,-\sbH+\sbY) -- (\sbH,-\sbH);
  \end{scope}}

% --- \tokencol{x}{y}{nbars}{width}{colour} : an N x C token set -----------
%     Deliberately NOT a slab: flat bars, visible gaps, no extrusion, so a
%     query set can never be mistaken for a feature map.
\newcommand{\tokencol}[5]{%
  \begin{scope}[shift={(#1,#2)}]
    \pgfmathsetmacro{\tkP}{0.098}
    \pgfmathsetmacro{\tkT}{0.056}
    \pgfmathsetmacro{\tkW}{0.5*#4}
    \pgfmathtruncatemacro{\tkN}{#3-1}
    \pgfmathsetmacro{\tkY}{0.5*(\tkN*\tkP+\tkT)}
    \foreach \tki in {0,...,\tkN} {
      \pgfmathsetmacro{\tkc}{\tkY-\tki*\tkP}
      \fill[#5!13] (-\tkW,\tkc-\tkT) rectangle (\tkW,\tkc);
      \draw[#5!75, line width=0.4pt] (-\tkW,\tkc-\tkT) rectangle (\tkW,\tkc);
    }
  \end{scope}}

% --- \maskstack{x}{y}{sidepx}{k}{colour}{zoom} : k planes of H x W --------
\newcommand{\maskstack}[6]{%
  \begin{scope}[shift={(#1,#2)}]
    \pgfmathsetmacro{\msH}{0.5*#6*\pxc{#3}}
    \pgfmathsetmacro{\msO}{0.30*\msH/2.5}
    \pgfmathtruncatemacro{\msN}{#4-1}
    \pgfmathsetmacro{\msS}{\msH/0.92}
    \foreach \msi in {\msN,...,0} {
      \pgfmathsetmacro{\msd}{\msi*\msO}
      \begin{scope}[shift={(\msd,\msd)}]
        \fill[white] (-\msH,-\msH) rectangle (\msH,\msH);
        \fill[#5!8]  (-\msH,-\msH) rectangle (\msH,\msH);
        \draw[#5!60, line width=0.35pt] (-\msH,-\msH) rectangle (\msH,\msH);
        \begin{scope}[scale=\msS]
          \fill[#5!35] \hull;
          \draw[#5!90, line width=0.3pt] \hull;
        \end{scope}
      \end{scope}}
  \end{scope}}

% --- \scorebar{x}{y}{h}{colour} : one scalar, as a short upright bar ------
\newcommand{\scorebar}[4]{%
  \begin{scope}[shift={(#1,#2)}]
    \draw[black!35, line width=0.35pt] (-0.10,0) rectangle (0.10,0.62);
    \fill[#4!55] (-0.10,0) rectangle (0.10,#3);
    \draw[#4!90, line width=0.4pt] (-0.10,0) rectangle (0.10,#3);
  \end{scope}}

% =========================================================================
% Reusable glyphs.  All are TikZ pics: \pic[scale=s] at (x,y) {name};
% =========================================================================

% --- lock: marks any component whose weights we never update -------------
\tikzset{lock/.pic={
  \fill[frozgray] (-0.105,-0.105) rectangle (0.105,0.055);
  \draw[frozgray, line width=0.65pt]
    (-0.060,0.055) -- (-0.060,0.105)
    arc[start angle=180, end angle=0, radius=0.060] -- (0.060,0.055);
  \fill[white] (0,-0.022) circle[radius=0.020];
}}

% --- planes: 3 offset stacked planes = a multi-scale feature pyramid -----
\tikzset{planes/.pic={
  \foreach \i/\o in {2/0.20, 1/0.10, 0/0.00} {
    \fill[datablue!\the\numexpr 10+8*\i\relax]
      (\o-0.24,\o-0.30) rectangle (\o+0.24,\o+0.30);
    \draw[datablue!70, line width=0.45pt]
      (\o-0.24,\o-0.30) rectangle (\o+0.24,\o+0.30);
  }
}}

% --- plate: a single feature plane (thin vertical plate, O4) -------------
\tikzset{plate/.pic={
  \fill[datablue!14] (-0.10,-0.34) rectangle (0.10,0.34);
  \draw[datablue!70, line width=0.45pt] (-0.10,-0.34) rectangle (0.10,0.34);
  \draw[datablue!45, line width=0.35pt] (-0.10,0.20) -- (0.10,0.20);
  \draw[datablue!45, line width=0.35pt] (-0.10,-0.20) -- (0.10,-0.20);
}}

% --- ckpt: an isometric slab = one stored checkpoint / weight vector -----
\tikzset{ckpt/.pic={
  \fill[frozgray!18] (-0.30,-0.13) rectangle (0.24,0.13);
  \fill[frozgray!32] (-0.30,0.13) -- (-0.24,0.21) -- (0.30,0.21) -- (0.24,0.13) -- cycle;
  \fill[frozgray!26] (0.24,-0.13) -- (0.30,-0.05) -- (0.30,0.21) -- (0.24,0.13) -- cycle;
  \draw[frozgray, line width=0.45pt]
    (-0.30,-0.13) rectangle (0.24,0.13);
  \draw[frozgray, line width=0.45pt]
    (-0.30,0.13) -- (-0.24,0.21) -- (0.30,0.21) -- (0.24,0.13)
    (0.30,0.21) -- (0.30,-0.05) -- (0.24,-0.13);
}}

% --- masked: pictogram of an output instance mask (hull on light field) --
\tikzset{maskicon/.pic={
  \fill[datablue!5] (-0.30,-0.30) rectangle (0.30,0.30);
  \draw[datablue!60, line width=0.5pt] (-0.30,-0.30) rectangle (0.30,0.30);
  \begin{scope}[rotate=28]
    \fill[datablue!45] (0,0.05) ellipse[x radius=0.235, y radius=0.062];
    \draw[datablue, line width=0.4pt] (0,0.05) ellipse[x radius=0.235, y radius=0.062];
  \end{scope}
  \begin{scope}[rotate=-12]
    \fill[datablue!45] (0.02,-0.19) ellipse[x radius=0.145, y radius=0.045];
    \draw[datablue, line width=0.4pt] (0.02,-0.19) ellipse[x radius=0.145, y radius=0.045];
  \end{scope}
}}

% =========================================================================
% \sarchip{x}{y}{side}{flip}  -- schematic SAR amplitude chip: dark speckled
% field with two bright elongated hulls.  flip = 0 identity, 1 horizontal
% mirror (so the TTA figure can show the geometry of the view, not a texture).
% =========================================================================
\newcommand{\sarchip}[4]{%
  \begin{scope}[shift={(#1,#2)}]
    \pgfmathsetmacro{\chSz}{#3}
    \pgfmathsetmacro{\chHw}{#3/2}
    \pgfmathsetmacro{\chSg}{1-2*(#4)}
    \pgfmathsetmacro{\chRa}{0.012*#3}
    \fill[sarink] (-\chHw,-\chHw) rectangle (\chHw,\chHw);
    \begin{scope}
      \clip (-\chHw,-\chHw) rectangle (\chHw,\chHw);
      % speckle
      \foreach \chK in {1,...,30} {
        \pgfmathsetmacro{\chX}{(rnd-0.5)*1.94*\chHw}
        \pgfmathsetmacro{\chY}{(rnd-0.5)*1.94*\chHw}
        \pgfmathsetmacro{\chO}{0.13+rnd*0.34}
        \fill[white, opacity=\chO] (\chX,\chY) circle[radius=\chRa];
      }
      % coastline / dock structure (bright straight clutter)
      \pgfmathsetmacro{\chDa}{\chSg*0.36*\chSz}
      \pgfmathsetmacro{\chDb}{\chSg*0.46*\chSz}
      \pgfmathsetmacro{\chDc}{\chSg*0.30*\chSz}
      \pgfmathsetmacro{\chDd}{\chSg*0.20*\chSz}
      \pgfmathsetmacro{\chDy}{0.5*\chSz}
      \fill[white, opacity=0.40]
        (\chDa,-\chDy) -- (\chDb,-\chDy) -- (\chDc,\chDy) -- (\chDd,\chDy) -- cycle;
      % two hulls
      \pgfmathsetmacro{\chAx}{\chSg*(-0.14)*\chSz}
      \pgfmathsetmacro{\chAy}{0.13*\chSz}
      \pgfmathsetmacro{\chAr}{\chSg*32}
      \pgfmathsetmacro{\chArx}{0.17*\chSz}
      \pgfmathsetmacro{\chAry}{0.048*\chSz}
      \fill[white, opacity=0.92, rotate around={\chAr:(\chAx,\chAy)}]
        (\chAx,\chAy) ellipse [x radius=\chArx, y radius=\chAry];
      \pgfmathsetmacro{\chBx}{\chSg*0.04*\chSz}
      \pgfmathsetmacro{\chBy}{-0.22*\chSz}
      \pgfmathsetmacro{\chBr}{\chSg*(-14)}
      \pgfmathsetmacro{\chBrx}{0.11*\chSz}
      \pgfmathsetmacro{\chBry}{0.036*\chSz}
      \fill[white, opacity=0.80, rotate around={\chBr:(\chBx,\chBy)}]
        (\chBx,\chBy) ellipse [x radius=\chBrx, y radius=\chBry];
    \end{scope}
    \draw[black!55, line width=0.4pt] (-\chHw,-\chHw) rectangle (\chHw,\chHw);
  \end{scope}}

% =========================================================================
% \hull{}  -- schematic ship-mask contour, ~1.3 cm wide, centred at origin.
% Used by the BPR and fusion figures so every ship outline in the paper is
% literally the same curve.
% =========================================================================
\newcommand{\hull}{plot[smooth cycle, tension=0.72] coordinates
  {(0.72,0.00) (0.34,0.26) (-0.22,0.33) (-0.64,0.17) (-0.62,-0.12)
   (-0.10,-0.30) (0.42,-0.22)}}
